\documentclass[11pt]{article}
\usepackage[T1]{fontenc}
\usepackage{lmodern,amsmath,amssymb,amsthm,microtype,xcolor}
\usepackage[letterpaper,margin=1in]{geometry}
\usepackage[colorlinks=true,linkcolor=blue!45!black,urlcolor=blue!45!black]{hyperref}
\newtheorem{theorem}{Theorem}[section]
\newtheorem{lemma}[theorem]{Lemma}
\newtheorem{proposition}[theorem]{Proposition}
\newtheorem{corollary}[theorem]{Corollary}
\newtheorem{remark}[theorem]{Remark}
\setlength{\emergencystretch}{3em}
\title{Structural reductions for the second-order cycle bound}
\author{Collaborative research working draft}
\date{11 November 2026}
\begin{document}
\maketitle
\begin{abstract}
The unrestricted conjectural error $O(r2^{r/2})$ remains open.
This report proves it for every fixed cyclic cut order and confines a sufficiently large minimal counterexample to a much smaller structural class.
Its girth $g$ and minimum cyclic cut order $k$ must satisfy
\[
 \frac12\log_2\log_2r-\frac{11}{2}\log_2\log_2\log_2r-O(1)
 <k=g
 <2\log_2\log_2r+2\log_2\log_2\log_2r+O(1),
\]
and every minimum cyclic cut must isolate a shortest cycle.
The minimal-counterexample statement explicitly assumes the sharp bound at all smaller ranks.

The main new tools are disconnected-region localization, quantitative terminal-path stability, and a uniform affine-forest estimate for high-girth poles.
If a minimum cut has two noncycle shores, the latter gives
\[
 c(G)\le\left(\frac14+O(g^{13/2}2^{-g/4})\right)2^r.
\]
At logarithmic girth, the established nested-chain argument works without any short-cycle deletions or lower bound on the positive forest excess.
Its two large shores force the same quarter-space conclusion and yield an exact edge-count dichotomy.
The report also gives explicit growing-cut and growing-girth ranges and records the limitations of several proposed shortcuts.
\end{abstract}

\section*{Status and scope}
The fixed-cut induction, the fixed-girth proof, quantitative localization, the high-girth terminal-forest theorem, and the stated refinements have passed four independent full mathematical audits.
The new theorems are not formalized in Lean.
Publication priority has not been established.
All constants may be very large.
The subscript in $O_K$ is essential: this report does not prove a constant uniform over unbounded cut orders.

The earlier six-cut classification remains a separate independent proof.
The induction below also covers both large shores of every fixed seven-cut, without solving the auxiliary terminal-state relaxation.
Numerical searches of that relaxation are not used here.

Rank means $|E|-|V|+\operatorname{comp}$.
Paths and cycles are simple, with parallel edges distinguished.
A cyclic cut has a cycle on each shore.
A normalized $k$-pole has $k$ distinct degree-two ports and all other degrees three.
Write $\psi(r)$ for the maximum cycle count at rank $r$.

\section{Inputs from the preceding research packets}
The following statements are inputs, not new claims of this report.
Their proofs appear in the accompanying path-stability, small-cut, and polynomial-error packets.

\begin{theorem}[Established cycle bounds]\label{thm:sharp-order-five-cut-error}
Uniformly over graphs of rank $r$, one has $c(G)\le(1/2+o(1))2^r$.
For connected loopless cubic multigraphs with a cyclic cut of size at most five,
\[
 c(G)\le2^{r-1}+C r2^{r/2}
\]
with an absolute constant $C$.
The unrestricted quantitative estimate is
\[
 \psi(r)\le2^{r-1}+O\!\left(2^r(\log_2r)^{11/2}/\sqrt r\right).
\]
\end{theorem}

\begin{proposition}[Established genuine five-cut deficit]\label{prop:five-cut-universal-deficit}
Let $G$ have a cyclic five-cut and no smaller cyclic cut.
Its shores are normalized five-poles of ranks $h\le l$.
Their crossing-cycle contribution, divided by $2^{h+l}$, is at most $15/2+o(1)$ as $l\to\infty$, uniformly in $h\ge1$.
Their internal-cycle contribution in the same normalization is at most $1/4+o(1)$.
Consequently $c(G)/2^{h+l}\le31/4+o(1)$.
\end{proposition}

The crossing bound uses five matching-augmentation inequalities and distinct representatives in the Petersen graph.
The internal bound follows from $c(X)\le2^h-1$ and $c(Y)\le(1/2+o(1))2^l$.
The full proof appears in \nolinkurl{small-cut-sharp-bound.pdf}.

\begin{theorem}[Established high-density path theorem]\label{thm:high-density-pair-small-cyclic-cut}
For $0<\varepsilon\le1/2$, let $C_\varepsilon$ be the additive terminal-path stability constant.
Set
\[
 B_{k,\varepsilon}=(k+1)\bigl(9\lceil C_\varepsilon\rceil+8k+9\bigr).
\]
A normalized $k$-pole of rank at least $B_{k,\varepsilon}$ with terminal-path density at least $1/2+\varepsilon$ contains two disjoint cyclic vertex sets.
They avoid all ports and have boundary at most four in every cubic completion.
\end{theorem}

\begin{theorem}[Established protected version]\label{thm:high-density-protected-cyclic-sets}
The preceding cyclic sets can avoid any fixed $m$ protected vertices if the rank is at least
\[
 (k+m+1)\bigl(9\lceil C_\varepsilon\rceil+8k+4m+9\bigr).
\]
\end{theorem}

The stability constant is supplied by \nolinkurl{terminal-path-stability.pdf}.
The high-density refinements are proved in \nolinkurl{small-cut-sharp-bound.pdf}.
The unrestricted estimate is proved in \nolinkurl{polynomial-cycle-error.pdf}.
The new quantitative bound on $C_\varepsilon$ below still falls short of the unrestricted sharp error.

\begin{theorem}[Established marked forest endpoint]
\label{thm:parallel-links-log-five-half-forest}
There is an absolute constant $M$ such that the following holds for all sufficiently large ranks $\rho$.
Let $H$ be a connected subcubic multigraph, and let $P$ be a nonempty connected marked vertex set containing every vertex of degree below three.
If
\[
 |P|\le\rho/(\log_2\rho)^2,
\]
then a uniform even subgraph of $H$ is a forest on $V(H)\setminus P$ with probability at most
\[
 \frac12+M\frac{(\log_2\rho)^{11/2}}{\sqrt\rho}.
\]
\end{theorem}

This is the connected-marked-set case of the endpoint theorem in \nolinkurl{polynomial-cycle-error.pdf}.
Its proof does not depend on additive path stability, so its use below is noncircular.

We also recall the precise part of the path-stability proof that will be reused.
A connected terminal network with path density above one half has a rank-preserving cubic expansion $\Gamma$.
Its terminals have degree one or two; all other vertices have degree three.
A shortest terminal path $A$ has length no greater than the original terminal distance.
Add a new terminal path of length two to obtain $H$, and mark its union $W$ with $A$.
The path injection bounds the original path density by the forest probability outside $W$.

For a threshold $\Theta\ge2^{36}$, call a set $R\subseteq V(\Gamma)\setminus V(A)$ good if
\[
 |R|\ge\Theta,\qquad |\delta_\Gamma R|\le |R|/(4\log_2^2|R|).
\]
The growing-window lemma says that, if no good set exists, each distance layer has fewer than $\Theta$ vertices outside $A$.
The additive layer argument then gives the constant
\[
 C_\varepsilon=\frac{2L(L+1)2^{L-1}}{\varepsilon},
 \qquad L=20\lceil\Theta\rceil+2.
\]
These are the cubic reduction, growing-window lemma, and additive estimate in the preceding terminal-path stability packet.
Only the old argument excluding good sets is replaced below.

\section{An induction over the cyclic cut order}
\label{sec:path-fixed-cut-induction}

Deleting one edge of a cyclic cut reduces its order by one after two suppressions.
The uniform terminal-path bound controls the cycles lost in this deletion.
This gives the sharp error for every fixed cyclic cut order.

Write $\lambda_c(G)$ for the minimum size of an edge cut whose two shores both contain cycles.
Set $\lambda_c(G)=\infty$ when no such cut exists.

\begin{lemma}[Smoothing one edge of a minimum cyclic cut]
\label{lem:path-minimum-cyclic-cut-smoothing}
Let $G$ be a connected cubic multigraph of girth at least five, and suppose $\lambda_c(G)=k\ge5$.
Choose a minimum cyclic cut $\delta_G A$, an edge $e=uv$ in it, and put $N=G-e$.
Suppress the degree-two vertices $u,v$ in $N$, obtaining $G'$.
Then
\begin{enumerate}
\item $G'$ is connected and loopless cubic, with $\beta(G')=\beta(G)-1$;
\item the two original shores give a cyclic $(k-1)$-cut in $G'$, with unchanged shore ranks;
\item $\lambda_c(G')=k-1$;
\item $c(G)=c(G')+P_{uv}(N)$.
\end{enumerate}
\end{lemma}

\begin{proof}
Each shore of the minimum cut is connected.
If one shore were disconnected, a cyclic component on that shore would have boundary smaller than $k$.
Its complement would still contain the opposite shore's cycle.

Each vertex has at most one incidence in the cut.
Otherwise its internal degree is at most one, so it lies on no cycle in its shore.
Moving it across preserves a cycle on each side and decreases the cut size.
Thus $u$ has two neighbors inside $A$, and $v$ has two inside its opposite shore.
The girth assumption makes each pair of internal neighbors distinct.

Every bridge in a cubic multigraph is a cyclic cut.
Indeed, a bridge shore with $n$ vertices and rank $h$ satisfies $n=2h-1$, so $h\ge1$.
The hypothesis therefore makes $G$ bridgeless, and $N$ is connected.
Deleting $e$ decreases rank by one.
Suppressing $u,v$ preserves rank and connectedness and creates no loop.
This proves the first assertion.
The suppression occurs wholly within each original shore.
It preserves both shore ranks, while the other $k-1$ cut edges remain unchanged.
This proves the second assertion.

To prove minimality, take any cyclic cut of size $j$ in $G'$.
Lift its shore partition when restoring each suppressed vertex.
If the endpoints of its smoothed edge lie on the same shore, put the restored vertex on that shore.
If they lie on opposite shores, put it on either shore.
The corresponding two-edge path then has the same number of cut edges as the smoothed edge.
Every cycle contained in either shore lifts to a cycle on that shore.
Finally restore $e$, which increases the boundary size by at most one.
This gives a cyclic cut of size at most $j+1$ in $G$.
Hence $j\ge k-1$, proving the third assertion.

Suppression gives a bijection between the cycles of $N$ and those of $G'$.
The cycles of $G$ containing $e$ correspond bijectively to the simple $u$--$v$ paths in $N$.
This proves the final identity.
\end{proof}

\begin{lemma}[Uniform path bound after one edge deletion]
\label{lem:path-cut-deletion-path-bound}
Let $G$ range over connected loopless cubic multigraphs of rank $r$ with no cyclic cut of size at most four.
For every edge $e=uv$, uniformly in $G$ and $e$,
\[
 P_{uv}(G-e)\le(1/2+o(1))2^{r-1}
 \qquad(r\to\infty).
\]
\end{lemma}

\begin{proof}
The graph $G$ is bridgeless by the bridge observation in Lemma~\ref{lem:path-minimum-cyclic-cut-smoothing}.
Thus $G-e$ is a connected normalized two-pole of rank $r-1$, with terminals $u,v$.
Its cubic completion obtained by restoring $e$ is $G$.
For each fixed $\varepsilon>0$, Theorem~\ref{thm:high-density-pair-small-cyclic-cut} excludes path density at least $1/2+\varepsilon$ once $r-1\ge B_{2,\varepsilon}$.
The threshold is independent of $G$ and $e$, which proves the uniform estimate.
\end{proof}

\begin{theorem}[A relative deficit at every fixed minimum cut order]
\label{thm:path-fixed-cyclic-order-deficit}
For every fixed integer $k\ge5$, connected loopless cubic multigraphs of rank $r$ with $\lambda_c(G)=k$ satisfy
\[
 c(G)\le\left(\frac12-2^{-k-1}+o_k(1)\right)2^r
 \qquad(r\to\infty),
\]
uniformly over these graphs.
\end{theorem}

\begin{proof}
The five-cut proof gives the required base case with its explicit relative gap.
Indeed, for $\lambda_c(G)=5$, write the shore ranks as $h\le l$.
Proposition~\ref{prop:five-cut-universal-deficit} and the internal-cycle estimate in Theorem~\ref{thm:sharp-order-five-cut-error} give
\[
 \frac{c(G)}{2^{h+l}}\le\frac{31}{4}+o(1).
\]
Since $r=h+l+4$, this is
\[
 \frac{c(G)}{2^r}\le\frac{31}{64}+o(1)
 =\frac12-2^{-6}+o(1).
\]

Suppose the theorem holds for $k-1\ge5$.
Let $\lambda_c(G)=k$ and $\beta(G)=r$.
At sufficiently large rank, $G$ has girth at least five.
Otherwise a cycle of length at most four has a cyclic complement and gives an excluded cut of size at most four.
Apply Lemma~\ref{lem:path-minimum-cyclic-cut-smoothing} to one edge of a minimum cyclic cut.
The resulting graph $G'$ has rank $r-1$ and minimum cyclic cut order $k-1$.
The induction hypothesis gives
\[
 c(G')\le\left(\frac12-2^{-k}+o_k(1)\right)2^{r-1}.
\]
Lemma~\ref{lem:path-cut-deletion-path-bound} gives
\[
 P_{uv}(G-e)\le(1/2+o(1))2^{r-1}.
\]
Adding these estimates through the exact cycle identity yields
\[
 c(G)\le\left(1-2^{-k}+o_k(1)\right)2^{r-1}
 =\left(\frac12-2^{-k-1}+o_k(1)\right)2^r.
\]
All errors are uniform at each fixed induction step.
This completes the induction.
\end{proof}

\begin{corollary}[Sharp error for every fixed cyclic cut order]
\label{cor:path-all-fixed-cyclic-cuts-sharp}
For every fixed integer $K\ge1$, there is a constant $C_K$ such that every connected loopless cubic multigraph $G$ of rank $r$ with a cyclic cut of size at most $K$ satisfies
\[
 c(G)\le2^{r-1}+C_K r2^{r/2}.
\]
For all sufficiently large ranks, if $G$ has no cyclic cut of size at most four, then the stronger uniform estimate holds:
\[
 c(G)\le\left(\frac12-2^{-K-1}+o_K(1)\right)2^r<2^{r-1}.
\]
The second assertion is vacuous when $K\le4$.
\end{corollary}

\begin{proof}
For a cyclic cut of size at most four, use the established sharp four-cut theorem.
Otherwise the minimum cyclic cut order is one of $5,\ldots,K$.
Apply Theorem~\ref{thm:path-fixed-cyclic-order-deficit} to these finitely many orders.
The ranks below the resulting thresholds are absorbed into $C_K$, using $c(G)\le2^r-1$.
\end{proof}

Consequently, any sequence of above-half cubic multigraphs with ranks tending to infinity either has small cyclic cuts or has minimum cyclic cut order tending to infinity, in the extended sense.
After the sharp four-cut class is removed, a counterexample sequence to the unrestricted sharp-order bound must have cyclic connectivity tending to infinity.
The unrestricted bound $O(r2^{r/2})$ remains open.

For sufficiently large cubic rank, the minimum cyclic cut order is finite and logarithmic in the rank.
Indeed, the breadth-first tree bound gives a shortest cycle of length
\[
 g\le2\log_2(2r-2)+2<r.
\]
A shortest cycle is induced at these ranks.
Deleting its vertices leaves rank at least $r-g>0$, by the cubic degree count.
It therefore defines a cyclic cut of size $g$.
Thus the remaining unbounded-order regime still satisfies
\[
 \lambda_c(G)\le2\log_2(2r-2)+2.
\]

\subsection{A finite recurrence for a growing cut order}

For an integer $s$, define the edge-deletion path excess
\[
 q(s)=\sup_{G,e}\left(2^{1-s}P_{uv}(G-e)-\frac12\right)_+,
\]
where $G$ has rank $s$, is connected loopless cubic, has no cyclic cut of size at most four, and $e=uv$.
Here and below, cubic graphs may have parallel edges.
Set a supremum over an empty class equal to zero.
Lemma~\ref{lem:path-cut-deletion-path-bound} gives $q(s)\to0$.
Define the nonnegative five-cut base error by
\[
 a_5(s)=\sup_{\substack{\beta(G)=s\\\lambda_c(G)=5}}
 \left(2^{-s}c(G)-\frac{31}{64}\right)_+.
\]
The uniform five-cut estimate gives $a_5(s)\to0$.

\begin{proposition}[Finite cut-order recurrence]
\label{prop:path-finite-cut-recurrence}
Let $G$ be connected loopless cubic, with rank $r$ and minimum cyclic cut order $k\ge5$.
Assume $r-k+5\ge10$.
Then
\begin{align*}
 \frac{c(G)}{2^r}
 \le{}&\frac12-2^{-k-1}
 +2^{-(k-5)}a_5(r-k+5)\\
 &+\sum_{i=0}^{k-6}2^{-i-1}q(r-i).
\end{align*}
The sum is empty when $k=5$.
\end{proposition}

\begin{proof}
Apply Lemma~\ref{lem:path-minimum-cyclic-cut-smoothing} repeatedly, each time to an edge of a minimum cyclic cut.
After $j=k-5$ steps, the rank is $r-j$ and the minimum cyclic cut order is five.
Every intermediate rank is at least ten, so the absence of a cyclic cut of size at most four guarantees the needed girth condition.
The exact cycle identity and the definition of $q$ give at step $i$
\[
 c(G_i)\le c(G_{i+1})+(1/2+q(r-i))2^{r-i-1}.
\]
Sum these identities for $0\le i<j$ and use
\[
 c(G_j)\le\left(\frac{31}{64}+a_5(r-j)\right)2^{r-j}.
\]
The geometric sum of the half terms gives the stated formula.
\end{proof}

For a sequence with $k=k(r)\ge5$, the same argument gives a strict half-bound deficit whenever
\[
 r-k\longrightarrow\infty,
 \qquad
 2^k\max_{r-k+6\le s\le r}q(s)\longrightarrow0.
\]
For $k=5$, the maximum condition is unnecessary.
The base-error term is then negligible relative to $2^{-k-1}$, as is the weighted path-excess sum.

There is also an explicit sufficient threshold in terms of additive path stability.
Choose $R_5$ so that $a_5(s)\le1/256$ for $s\ge R_5$.
Set $\varepsilon=2^{-k-3}$ and
\[
 B_{2,\varepsilon}=3\bigl(9\lceil C_\varepsilon\rceil+25\bigr).
\]
If
\[
 r-k+5\ge\max\{10,R_5,B_{2,2^{-k-3}}\},
\]
then the path theorem gives $q(s)\le2^{-k-3}$ at every rank in the recurrence.
The two error terms are each at most one quarter of $2^{-k-1}$.
Consequently
\[
 c(G)\le\left(\frac12-2^{-k-2}\right)2^r.
\]
This condition permits a cut order that grows sufficiently slowly with rank.
Its quantitative limitation is the dependence of $C_\varepsilon$ on $\varepsilon$.

\section{A strict deficit at bounded girth}
\label{sec:path-fixed-girth}

A shortest cycle gives a sequence of terminal-path bounds.
The argument permits a disconnected complement of that cycle.

The proof uses the asymptotic cycle bound and the protected-vertices theorem.
The latter is Theorem~\ref{thm:high-density-protected-cyclic-sets}.
Thus its quantitative constants ultimately depend on additive terminal-path stability.

\begin{lemma}[Prefix deletion along a shortest cycle]
\label{lem:path-shortest-cycle-prefix}
Let $G$ be a connected loopless cubic multigraph of rank $r$ and girth $g\ge5$.
Let $C=v_0v_1\cdots v_{g-1}v_0$ be a shortest cycle.
Write $w_i$ for the unique neighbor of $v_i$ outside $C$.
Put $N=G-v_0v_{g-1}$ and $b=\beta(N)=r-1$.
For $1\le p\le g-1$, delete $v_0,\ldots,v_{p-1}$ from $N$.
Let $K_p$ be the component containing $v_{g-1}$.
Then
\[
 \beta(K_p)\le b-p.
\]
If $p\le g-2$, then $K_p$ is a normalized pole with at most $g$ ports.
If $p=g-1$, deleting the pendant vertex $v_{g-1}$ from $K_p$ gives such a pole of the same rank.
Whenever $w_{p-1}\in K_p$, it is a port of the resulting pole.
Its terminal partner is $v_{g-1}$ for $p\le g-2$, and $w_{g-1}$ for $p=g-1$.
These terminal partners are distinct from $w_{p-1}$.
\end{lemma}

\begin{proof}
The cycle $C$ is induced.
An outside vertex cannot have two neighbors on $C$.
Otherwise the shorter arc between those neighbors gives a cycle of length at most
\[
 \left\lfloor\frac g2\right\rfloor+2<g.
\]
Consequently all $w_i$ are distinct, and every component of $G-V(C)$ has minimum degree at least two.
Each such component contains a cycle.

Deleting the prefix removes $p$ cycle-path edges and $p$ outside edges.
If the remaining graph has $c_p$ components, its total rank is
\[
 b-p+c_p-1.
\]
Every component other than $K_p$ avoids the remaining cycle path.
It is therefore an entire component of $G-V(C)$ and has rank at least one.
Subtracting these $c_p-1$ ranks gives $\beta(K_p)\le b-p$.

For $p\le g-2$, the two endpoints of the remaining cycle path have degree two.
Its other vertices have degree three.
An outside vertex loses at most one incident edge, so its degree remains at least two.
The degree-two vertices are these two endpoints and some of $w_0,\ldots,w_{p-1}$.
There are at most $p+2\le g$ such vertices.
In particular, $w_{p-1}$ has degree two whenever it belongs to $K_p$.

For $p=g-1$, the remaining cycle vertex has only its outside edge.
Delete that pendant vertex and move the target to $w_{g-1}$.
The graph becomes a component of $G-V(C)$ and its rank does not change.
Its ports are some of $w_0,\ldots,w_{g-1}$.
The distinctness of the outside neighbors proves the final assertion.
\end{proof}

\begin{theorem}[Fixed-girth deficit]
\label{thm:path-fixed-girth-deficit}
Fix an integer $g\ge5$.
Let $G$ range over connected loopless cubic multigraphs of rank $r$, girth $g$, and no cyclic cut of size at most four.
Then, uniformly over these graphs,
\[
 c(G)\le 2^{r-1}-2^{r-g-1}+o_g(2^r)
 \qquad(r\to\infty).
\]
In particular, $c(G)<2^{r-1}$ for all sufficiently large $r$.
\end{theorem}

\begin{proof}
Use the notation of Lemma~\ref{lem:path-shortest-cycle-prefix}.
Let $P$ count the simple $v_0$--$v_{g-1}$ paths in $N$.
The cycle path $v_0v_1\cdots v_{g-1}$ contributes one.
Every other path first leaves this fixed path through $v_iw_i$, for some $0\le i\le g-2$.
Its remaining segment avoids $v_0,\ldots,v_i$.

Put $p=i+1$.
If $w_i\notin K_p$, the corresponding branch contributes zero.
Otherwise its contribution is exactly the terminal-path count in the pole from Lemma~\ref{lem:path-shortest-cycle-prefix}.
Write $q_p$ for that pole's rank.
Then $q_p\le b-p$.

Fix $\varepsilon>0$.
There is a threshold $B=B(g,\varepsilon)$ such that every nonzero branch of rank at least $B$ has terminal-path count at most
\[
 (1/2+\varepsilon)2^{q_p}.
\]
Indeed, apply the protected-vertices theorem for each possible number of ports, between two and $g$.
Protect all remaining vertices of $C$.
The theorem supplies two disjoint cyclic vertex sets if this bound fails.
They avoid every port and every protected vertex.
Every changed incidence in passing back to $G$ is at a port.
Both sets therefore retain their cycles and boundaries of size at most four in $G$, which is impossible.

A branch of rank less than $B$ has at most $2^B$ terminal paths.
This follows from the usual injection of terminal paths into their affine parity class.
Hence every branch, uniformly in its rank, has at most
\[
 (1/2+\varepsilon)2^{b-p}+2^B
\]
paths.
Summing over the $g-1$ branches gives
\begin{align*}
 P
 &\le 1+(1/2+\varepsilon)\sum_{p=1}^{g-1}2^{b-p}+(g-1)2^B\\
 &=1+(1/2+\varepsilon)\bigl(2^b-2^{b-g+1}\bigr)+(g-1)2^B.
\end{align*}
First let $r\to\infty$, and then let $\varepsilon\downarrow0$.
The threshold $B$ depends only on $g$ and $\varepsilon$.
Therefore
\[
 P\le 2^{b-1}-2^{b-g}+o_g(2^b).
\]

Every cycle of $G$ either avoids $v_0v_{g-1}$ or corresponds to one of these terminal paths.
Thus $c(G)=c(N)+P$.
The asymptotic cycle bound applies to the connected graph $N$ of rank $b$ and gives
\[
 c(N)\le2^{b-1}+o(2^b).
\]
Adding the last two estimates and substituting $b=r-1$ proves the theorem.
\end{proof}

\begin{corollary}[The remaining girth must tend to infinity]
\label{cor:path-unbounded-girth-reduction}
Fix an integer $g_0\ge2$.
There is a rank threshold $R(g_0)$ such that every connected loopless cubic multigraph $G$ with
\[
 \beta(G)\ge R(g_0),\qquad \operatorname{girth}(G)\le g_0,
\]
and no cyclic cut of size at most four satisfies $c(G)<2^{\beta(G)-1}$.
Consequently, any sequence with ranks tending to infinity, no such cyclic cuts, and cycle counts above the half bound has girth tending to infinity.
\end{corollary}

\begin{proof}
If a cycle has at most four vertices, deleting its vertices reduces rank by at most eight.
For rank greater than eight, its complement therefore contains a cycle.
Its boundary has at most four edges, which contradicts the cut assumption.
Thus the girth is at least five at sufficiently large rank.
Apply Theorem~\ref{thm:path-fixed-girth-deficit} to the finitely many values $5\le g\le g_0$.
\end{proof}

The argument does not require $G-V(C)$ to be connected.
It also does not require the cut around $C$ to have minimum cyclic order.
The rank bound in Lemma~\ref{lem:path-shortest-cycle-prefix} is the needed substitute for those assumptions.

\begin{corollary}[Every fixed-order cut with a bounded-rank shore]
\label{cor:path-bounded-shore-any-order}
Fix integers $k,H\ge1$.
For all sufficiently large ranks, the following statement holds.
Let $G$ be a connected loopless cubic multigraph with no cyclic cut of size at most four.
If $G$ has a cyclic cut of size at most $k$ with one shore of total rank at most $H$, then
\[
 c(G)<2^{\beta(G)-1}.
\]
Neither shore is required to be connected, and distinct cut edges may have a common endpoint.
\end{corollary}

\begin{proof}
Choose a cyclic component $A$ of the bounded-rank shore.
Write $h_A=\beta(G[A])\le H$ and $d_A=|\delta_G A|\le k$.
The cubic degree identity gives
\[
 |A|=2h_A+d_A-2\le2H+k-2.
\]
Thus $A$ contains a cycle of length at most $2H+k-2$.
Apply Corollary~\ref{cor:path-unbounded-girth-reduction} with the fixed bound $\max\{2,2H+k-2\}$.
\end{proof}

In particular, the remaining case for any fixed cut order has both shore ranks tending to infinity.
Combining this corollary with the established sharp bound for cyclic cuts of size at most four also gives
\[
 c(G)\le2^{r-1}+O_{k,H}(r2^{r/2})
\]
for every graph with such a bounded-rank cyclic shore, without the small-cut exclusion.
The finitely many ranks below the strict-deficit threshold are absorbed into the constant.

\subsection{A finite inequality and a growing-girth condition}

The branch poles already contain every changed incidence at a port.
Thus the unprotected small-cut theorem gives a slightly smaller explicit threshold.

\begin{proposition}[Finite shortest-cycle estimate]
\label{prop:path-fixed-girth-finite}
Let $G$ satisfy the hypotheses of Theorem~\ref{thm:path-fixed-girth-deficit}, with rank $r$ and girth $g\ge5$.
The value of $g$ need not be fixed here.
For $0<\varepsilon\le1/2$, put
\[
 B(g,\varepsilon)
  =(g+1)\bigl(9\lceil C_\varepsilon\rceil+8g+9\bigr),
 \qquad
 E(s)=\max\{\psi(s)-2^{s-1},0\}.
\]
Then
\begin{align*}
 c(G)\le{}&2^{r-1}-2^{r-g-1}
       +\varepsilon\bigl(2^{r-1}-2^{r-g}\bigr)
       +E(r-1)\\
       &\hspace{20mm}+1+(g-1)2^{B(g,\varepsilon)}.
\end{align*}
\end{proposition}

\begin{proof}
Each nonzero branch is a normalized pole with between two and $g$ ports.
Every incidence added in returning to $G$ is at one of its ports.
Theorem~\ref{thm:high-density-pair-small-cyclic-cut} therefore bounds each branch of rank at least $B(g,\varepsilon)$ by $(1/2+\varepsilon)2^{q_p}$.
A branch below this threshold has at most $2^{B(g,\varepsilon)}$ paths.
Use $q_p\le r-1-p$, sum the geometric series, and add
\[
 c(G-v_0v_{g-1})\le2^{r-2}+E(r-1).
\]
This gives the displayed finite inequality.
\end{proof}

For example, consider a sequence with ranks $r\to\infty$, girths $g=g(r)\ge5$, and no cyclic cut of size at most four.
Set $\varepsilon=2^{-g-2}$.
The sequence is eventually below the half bound if
\begin{equation}\label{eq:path-growing-girth-conditions}
 \frac{2^g E(r-1)}{2^r}\longrightarrow0,
 \qquad
 B(g,2^{-g-2})+g+\log_2 g-r\longrightarrow-\infty.
\end{equation}
Indeed, the $\varepsilon$ term is at most one quarter of the negative term $2^{r-g-1}$.
Both remaining error terms are negligible compared with that negative term under these conditions.

The proved endpoint cycle bound makes the first condition hold whenever
\[
 \frac{2^g(\log r)^{11/2}}{\sqrt r}\longrightarrow0.
\]
The second condition still depends on the additive stability constant $C_{2^{-g-2}}$.
No polynomial estimate for that constant is asserted here.

\section{Disconnected localization and quantitative path stability}
\label{sec:path-quantitative-localization}

Localization does not require the retained vertex set to be connected.
Removing that requirement avoids the component-counting step in the original stability proof.
The endpoint forest estimate then gives a quantitative additive stability constant.
All logarithms in this section have base two.

\begin{lemma}[Localization on a disconnected shore]
\label{lem:path-disconnected-localization}
Let $G$ be a connected multigraph, and let $R$ be a nonempty proper vertex set such that $G-R$ is connected.
Write
\[
 d=|\delta_G R|\ge1,\qquad c_R=\operatorname{comp}(G[R]).
\]
For each boundary edge, add a distinct new vertex and a spoke to its endpoint in $R$.
Join the new vertices into a cycle $W_R$, padding it to length three when $d<3$.
Let $G_R$ consist of $G[R]$, these spokes, and $W_R$.
Then $G_R$ is connected and
\[
 |V(W_R)|=\max\{d,3\},
 \qquad
 \beta(G_R)=\beta(G[R])+d-c_R+1.
\]
For uniform even subgraphs,
\[
 \Pr\bigl(\xi_G\cap E(G[R])\text{ is a forest}\bigr)
 =\Pr\bigl(\xi_{G_R}\cap E(G_R-V(W_R))\text{ is a forest}\bigr).
\]
If every vertex of $R$ has degree three in $G$, then
\[
 \beta(G_R)=\frac{|R|+d}{2}+1.
\]
\end{lemma}

\begin{proof}
Every component of $G[R]$ has a boundary edge, since $G$ is connected.
The new cycle therefore connects all these components.
Without padding, the construction adds $d$ vertices and $2d$ edges.
Its rank is $|E(G[R])|+d-|R|+1$, giving the stated formula.
Padding adds equally many edges and vertices and preserves this rank.

Contract all edges of the connected graph $G-R$ in $G$.
Separately, contract all edges of $W_R$ in $G_R$.
The resulting multigraphs agree, including their restrictions to $E(G[R])$.
Projection of a uniform cycle-space vector through a contraction is uniform on the image cycle space.
For surjectivity, complete each prescribed even boundary by a join in a spanning tree of the contracted component.
Thus the restrictions to $E(G[R])$ have the same law in both constructions.
This proves the probability identity.

Under the cubic assumption, $3|R|=2|E(G[R])|+d$.
Substitution into the rank formula gives $(|R|+d)/2+1$.
No connectedness assumption on $R$ was used.
\end{proof}

\begin{lemma}[A quantitative localized forest threshold]
\label{lem:path-polynomial-local-threshold}
There is an absolute constant $K_0$ with the following property.
For $0<\delta\le1/4$, one can choose an integer $R_\delta\ge2^{100}$ satisfying
\[
 R_\delta\le K_0\delta^{-2}\bigl(\log(2/\delta)\bigr)^{11}.
\]
Whenever a connected subcubic multigraph has rank $\rho\ge R_\delta$ and a connected marked set $P$ containing every vertex of degree below three, with
\[
 |P|\le\frac{\rho}{(\log\rho)^2},
\]
the probability that its uniform even subgraph is a forest outside $P$ is at most $1/2+\delta$.
\end{lemma}

\begin{proof}
Use the absolute constant $M$ and absolute rank threshold in Theorem~\ref{thm:parallel-links-log-five-half-forest}.
Increase that threshold to at least $2^{100}$.
At these ranks, a connected marked set satisfies its component bound.
It is enough to require
\[
 M\frac{(\log\rho)^{11/2}}{\sqrt\rho}\le\delta.
\]
The left side decreases for all sufficiently large $\rho$.
Choosing $\rho$ above a sufficiently large constant times
\[
 \delta^{-2}\bigl(\log(2/\delta)\bigr)^{11}
\]
ensures this inequality.
Indeed, the logarithm of this expression is $O(\log(2/\delta))$.
Choose the absolute constant large enough to cover both the endpoint threshold and the remaining compact range of $\delta$.
Rounding upward gives the asserted integer $R_\delta$.
\end{proof}

\begin{lemma}[Every sufficiently large good set is excluded]
\label{lem:path-direct-good-set}
Use the cubic expansion $\Gamma$, shortest terminal path $A$, graph $H$, and marked cycle $W$ from the general stability proof.
Let $\delta=\varepsilon/2$ and
\[
 \Theta=\max\{2^{36},2R_\delta\}.
\]
If the terminal-path density is at least $1/2+\varepsilon$, no set $R\subseteq V(\Gamma)\setminus V(A)$ can satisfy
\[
 |R|\ge\Theta,
 \qquad
 |\delta_\Gamma R|\le\frac{|R|}{4(\log|R|)^2}.
\]
\end{lemma}

\begin{proof}
Suppose such an $R$ exists, and put $m=|R|$.
Let $C_0$ be the component of $H-R$ containing $W$.
Set $R^+=V(H)\setminus V(C_0)$.
Then $R^+\supseteq R$, its complement is connected, and it avoids $V(W)$.
It need not be connected.
Every boundary edge of $R^+$ already belongs to $\delta_H R$, because distinct components of $H-R$ have no edge between them.
Hence, writing $d^+=|\delta_HR^+|$,
\[
 d^+\le|\delta_HR|=|\delta_\Gamma R|
 \le\frac{m}{4(\log m)^2}.
\]
All vertices of $R^+$ are cubic in $H$.

Apply Lemma~\ref{lem:path-disconnected-localization} to $R^+$.
The resulting graph has rank
\[
 \rho=\frac{|R^+|+d^+}{2}+1\ge m/2\ge R_\delta.
\]
It is subcubic, and every vertex outside its marked cycle is cubic.
The marked cycle is connected and has at most $\rho/(\log\rho)^2$ vertices.
For $d^+\ge3$, this follows from the increasing function $\Phi(t)=t/(\log t)^2$ on $t\ge8$ and
\[
 d^+\le\Phi(m)/4\le\Phi(m/2)\le\Phi(\rho).
\]
For $d^+<3$, padding adds only three marked vertices, which also meets that bound.

The terminal-path injection bounds $1/2+\varepsilon$ by the forest probability outside $W$ in $H$.
That event implies the forest event on $H[R^+]$.
The localization identity and Lemma~\ref{lem:path-polynomial-local-threshold} bound its probability by $1/2+\delta$.
This contradicts $\delta=\varepsilon/2$.
\end{proof}

\begin{theorem}[A single-exponential additive stability constant]
\label{thm:path-quantitative-additive}
There is an absolute constant $K_1$ such that the following holds.
Let $N$ be a connected graph of rank $b$ with distinct terminals $u,v$, and suppose
\[
 P_{uv}(N)\ge(1/2+\varepsilon)2^b,
 \qquad 0<\varepsilon\le1/2.
\]
Then every terminal path has length at least $b-C_\varepsilon$, where one may choose
\[
 C_\varepsilon\le
 2^{K_1\varepsilon^{-2}(\log(2/\varepsilon))^{11}}.
\]
The cubic expansion in the stability proof has distance layers of size at most
\[
 K_1\varepsilon^{-2}(\log(2/\varepsilon))^{11}+1.
\]
\end{theorem}

\begin{proof}
Lemma~\ref{lem:path-direct-good-set} replaces the original good-set dichotomy.
The unchanged growing-window lemma gives fewer than $\Theta$ unmarked vertices in each distance layer.
Lemma~\ref{lem:path-polynomial-local-threshold} bounds
\[
 \Theta\le K_2\varepsilon^{-2}(\log(2/\varepsilon))^{11}
\]
for an absolute $K_2$.
There is at most one marked path vertex in each layer.
This proves the layer-width statement after increasing the final constant.

Put $B=\lceil\Theta\rceil$ and $L=20B+2$.
The minimum-span window argument and greedy deletion of whole windows remain unchanged.
The already proved additive estimate therefore permits
\[
 C_\varepsilon=\frac{2L(L+1)2^{L-1}}{\varepsilon}.
\]
Taking logarithms, and using the displayed polynomial bound on $\Theta$, gives
\[
 \log C_\varepsilon
 \le K_1\varepsilon^{-2}(\log(2/\varepsilon))^{11}
\]
after increasing $K_1$.
The terminal reduction and rank-distance comparison then give the asserted path bound.
\end{proof}

\begin{corollary}[A quantitative edge-deletion path excess]
\label{cor:path-logarithmic-deletion-excess}
The function $q(s)$ from Proposition~\ref{prop:path-finite-cut-recurrence} satisfies
\[
 q(s)=O\!\left(\frac{(\log\log s)^{11/2}}{\sqrt{\log s}}\right).
\]
\end{corollary}

\begin{proof}
Put $x=\log s$, $z=\log x$, and choose a fixed $A\ge\max\{1,2\sqrt{K_1}\}$.
Set $\varepsilon=A z^{11/2}/\sqrt x$.
For sufficiently large $s$, this lies in $(0,1/2]$ and $\log(2/\varepsilon)\le z$.
Theorem~\ref{thm:path-quantitative-additive} then gives
\[
 C_\varepsilon
 \le2^{K_1\varepsilon^{-2}(\log(2/\varepsilon))^{11}}
 \le2^{x/4}=s^{1/4}.
\]
Consequently
\[
 B_{2,\varepsilon}=3(9\lceil C_\varepsilon\rceil+25)<s-1
\]
for large $s$.
The two-pole high-density small-cut theorem gives $q(s)\le\varepsilon$.
This proves the stated rate.
\end{proof}

\begin{corollary}[Growing cyclic cut orders]
\label{cor:path-loglog-growing-cut}
Let connected loopless cubic multigraphs $G$ have ranks $r\to\infty$ and minimum cyclic cut orders $k=k(r)\ge5$.
If
\[
 \frac{2^k(\log\log r)^{11/2}}{\sqrt{\log r}}\longrightarrow0,
\]
then
\[
 c(G)\le2^{r-1}-(1-o(1))2^{r-k-1}<2^{r-1}.
\]
In particular, this holds whenever
\[
 k\le(1/2-\eta)\log\log r
\]
for any fixed $\eta>0$.
\end{corollary}

\begin{proof}
The assumed condition gives $k=O(\log\log r)$, so $r-k\to\infty$.
Corollary~\ref{cor:path-logarithmic-deletion-excess} holds uniformly over the ranks from $r-k+6$ through $r$.
Its bound changes by a factor $1+o(1)$ across this interval.
Apply Proposition~\ref{prop:path-finite-cut-recurrence} and compare its error terms with $2^{-k-1}$.
Both are negligible under the displayed assumption.
For the final assertion, the ratio is at most
\[
 \frac{(\log\log r)^{11/2}}{(\log r)^\eta},
\]
which tends to zero.
\end{proof}

\begin{corollary}[The endpoint growing-order range]
\label{cor:path-endpoint-loglog-cut}
There is an absolute constant $C$ such that, for all sufficiently large $r$, every connected loopless cubic multigraph of rank $r$ with
\[
 5\le\lambda_c(G)=k
 \le\frac12\log\log r-\frac{11}{2}\log\log\log r-C
\]
satisfies
\[
 c(G)\le2^{r-1}-2^{r-k-2}.
\]
\end{corollary}

\begin{proof}
Choose an absolute $A_0$ such that the path-excess sum in Proposition~\ref{prop:path-finite-cut-recurrence} is at most
\[
 A_0\frac{(\log\log r)^{11/2}}{\sqrt{\log r}}
\]
throughout $k=O(\log\log r)$ at sufficiently large $r$.
This follows from Corollary~\ref{cor:path-logarithmic-deletion-excess} and the sum of its geometric weights, which is less than one.
Take $C$ so large that $A_0 2^{-C}\le1/8$.
The displayed cut-order bound then makes the path-excess sum at most $2^{-k-3}$.
For sufficiently large $r$, the base error satisfies $a_5(r-k+5)\le1/256$ uniformly in this range.
Its contribution is also at most $2^{-k-3}$.
Subtract these two errors from the deficit $2^{-k-1}$ in Proposition~\ref{prop:path-finite-cut-recurrence}.
The remaining deficit is at least $2^{-k-2}$, as asserted.
\end{proof}

The new localization step removes an unnecessary exponential loss in the stability constants.
It does not give the unrestricted half-exponent cycle error.
The remaining cyclic connectivity can still grow between the displayed double-logarithmic range and the general logarithmic upper bound.

\section{High girth shortens terminal geodesics}
\label{sec:four-high-connectivity-moore-paths}

All logarithms in this section have base two.
A normalized $k$-pole is connected, has $k$ vertices of degree two, and has degree three at every other vertex.
The first bound uses no cyclic-connectivity assumption.

\begin{lemma}[Moore balls along a geodesic]
\label{lem:four-pole-moore-geodesic}
Let $N$ be a normalized $k$-pole of order $n$ and girth at least an integer $g\ge3$.
Put
\[
 t=\left\lfloor\frac{g-1}{2}\right\rfloor,
 \qquad B_t=3\cdot2^t-2.
\]
Every shortest path in $N$, of length $\ell$, satisfies
\begin{equation}\label{eq:four-moore-geodesic-exact}
 \ell+1<(2t+1)\left(k+1+\frac{n}{B_t}\right)
 \le g\left(k+1+n2^{-g/2}\right).
\end{equation}
\end{lemma}

\begin{proof}
Write the path as $v_0v_1\cdots v_\ell$.
Choose centers at positions $t+j(2t+1)$ whenever the full corresponding interval lies in $[0,\ell]$.
There are
\[
 q=\left\lfloor\frac{\ell+1}{2t+1}\right\rfloor
\]
centers.
The radius-$t$ vertex balls around these centers are disjoint.
Otherwise two centers would have distance at most $2t$, contradicting the geodesic property.
At most $k$ balls contain a degree-two vertex.

Every remaining ball contains at least $B_t$ vertices.
Indeed, breadth-first paths of length at most $t$ cannot collide, since a collision gives a cycle of length at most $2t<g$.
All vertices at distance below $t$ have degree three.
The breadth-first tree therefore has $1+3(1+2+\cdots+2^{t-1})=B_t$ vertices.
Edges within the last sphere do not affect this count.
Thus $(q-k)_+B_t\le n$, and hence $q\le k+n/B_t$.
The floor defining $q$ gives the first bound.

Since $t\ge1$, one has $B_t\ge2^{t+1}\ge2^{g/2}$.
Also $2t+1\le g$.
These inequalities give the final bound.
\end{proof}

\begin{corollary}[Short cycles through every edge at high girth]
\label{cor:four-high-girth-edge-cycle}
Let $G$ be a connected bridgeless cubic graph of rank $r$ and girth at least $g\ge3$.
Every edge lies on a cycle of length less than
\[
 g\left(3+(2r-2)2^{-g/2}\right).
\]
In particular, for every fixed $\eta>0$, girth at least $(4+\eta)\log\log r$ gives cycles of length $o(r/\log^2r)$ through every edge.
\end{corollary}

\begin{proof}
For an edge $e=uv$, the graph $N=G-e$ is a connected normalized two-pole of order $2r-2$.
Its girth is at least that of $G$.
A shortest terminal path together with $e$ gives the desired cycle.
Apply Lemma~\ref{lem:four-pole-moore-geodesic} with $k=2$.

For the asymptotic assertion, use the integer lower bound $g_0=\lfloor(4+\eta)\log\log r\rfloor$ in that lemma.
Then $g_0=O(\log\log r)$ and $2^{-g_0/2}=O((\log r)^{-2-\eta/2})$.
Both terms in the asserted length bound are $o(r/\log^2r)$.
\end{proof}

\begin{corollary}[Endpoint forest error for high-girth edge deletion]
\label{cor:four-high-girth-deletion-excess}
There is an absolute constant $C$ with the following property.
Let $G$ be connected and bridgeless cubic, of sufficiently large rank $r$.
Write $x=\log r$ and $z=\log x$.
If
\[
 \operatorname{girth}(G)\ge
 \left\lceil4z+2\log z+12\right\rceil,
\]
then, for every edge $e=uv$,
\[
 \frac{P_{uv}(G-e)}{2^{r-1}}
 \le\frac12+C\frac{(\log r)^{11/2}}{\sqrt r}.
\]
The same conclusion holds under girth at least $(4+\eta)\log\log r$ for any fixed $\eta>0$ and sufficiently large $r$.
\end{corollary}

\begin{proof}
Use $g_0=\lceil4z+2\log z+12\rceil$ in Corollary~\ref{cor:four-high-girth-edge-cycle}.
For sufficiently large $z$, one has $g_0\le8z$ and
\[
 2^{-g_0/2}\le\frac{1}{64x^2z}.
\]
Thus every edge $e=uv$ lies on a cycle $W=A\cup\{e\}$ with
\[
 |V(W)|<24z+\frac{r}{4x^2}\le\frac{r}{x^2}.
\]
Here $A$ is a terminal path in $G-e$.
Each terminal path $P$ determines a distinct outside part $P\setminus E(A)$.
Distinctness follows because their symmetric difference cannot be a nonempty even set supported on the path $A$.
There are $2^{r-1}$ possible outside parts, and each path gives a forest outside $W$.

The outside-part law is the restriction of the uniform cycle-space law of $G$.
Being a forest outside the edges of $W$ implies being a forest outside its vertices.
Apply Theorem~\ref{thm:parallel-links-log-five-half-forest} with the connected marked set $V(W)$.
Every vertex of $G$ is cubic, and the displayed size bound supplies its remaining hypothesis.
This proves the estimate uniformly over all edges.
\end{proof}

\begin{proposition}[A rank-independent girth bound for path excess]
\label{prop:four-girth-path-excess}
There is an absolute constant $C$ such that every normalized two-pole $N$ of rank $b$ and girth at least $g\ge3$ satisfies
\[
 \left(\frac{P_{uv}(N)}{2^b}-\frac12\right)_+
 \le C g^{11/2}2^{-g/4}.
\]
\end{proposition}

\begin{proof}
Let $t=\lfloor(g-1)/2\rfloor$, and let $L_t$ be the distance-$t$ sphere from the terminal $u$.
The breadth-first expansion starts with two edges at $u$.
Until radius $t$, distinct expansion paths cannot reach the same vertex, because $2t<g$.
Only the other terminal $v$ can remove an expected child.
If its distance from $u$ is $s<t$, that removal costs $2^{t-s-1}$ vertices in the final sphere.
Otherwise there is no removal.
Consequently
\[
 |L_t|\ge2^{t-1}\ge2^{g/2-2}.
\]

Suppose the normalized path count is at least $1/2+\varepsilon$.
The quantitative localization proof applies directly to $N$, with no cubic expansion needed.
Its distance-layer bound, from Theorem~\ref{thm:path-quantitative-additive}, gives
\[
 |L_t|\le K_1\varepsilon^{-2}\bigl(\log(2/\varepsilon)\bigr)^{11}+1.
\]
Choose $C\ge4\sqrt{K_1+1}$ and put $\varepsilon=Cg^{11/2}2^{-g/4}$.
For sufficiently large $g$, this lies in $(0,1/2]$ and $\log(2/\varepsilon)\le g$.
The last upper bound is then at most
\[
 \frac1{16}2^{g/2}+1<2^{g/2-2},
\]
a contradiction.
Increase $C$ to cover the finitely many remaining values of $g$.
\end{proof}

The independent geometric checker is \texttt{four-high-connectivity-moore-check.py}.
Its output, \texttt{four-high-connectivity-moore-checks.json}, records 49 normalized pole fixtures.
They include edge-deleted and subdivided Petersen, Heawood, and honeycomb graphs, with girth up to twelve.
All 35,568 geodesic checks and 40 terminal-sphere checks pass.
These finite checks validate the counting implementation; the proof above supplies the general statement.

\section{A rank dichotomy for high-girth shores}
\label{sec:four-high-connectivity-entropy-moore}

The following average-growth bound uses a stationary non-backtracking walk.
It avoids an external irregular Moore-bound theorem.
All logarithms and entropies have base two.

\begin{lemma}[An entropy Moore bound for normalized poles]
\label{lem:four-entropy-moore}
Let $N$ be a normalized $k$-pole of rank $h\ge2$, order $n$, and girth $g\ge3$.
Then
\[
 n=2h+k-2,
 \qquad
 g\le4+2\left(1+\frac{k}{3(h-1)}\right)\log n.
\]
\end{lemma}

\begin{proof}
The degree sum gives $2h-2$ cubic vertices and $k$ degree-two vertices.
Hence $n=2h+k-2$ and $m=|E(N)|=3h+k-3$.

Start with a uniformly chosen directed edge.
At each step, choose uniformly among edges that do not reverse the preceding step.
The directed-edge distribution remains uniform.
Indeed, a directed edge leaving $v$ has $\deg(v)-1$ possible predecessors, each assigning it probability $1/(\deg(v)-1)$.
Thus the conditional entropy of each further step is
\[
 \sigma=\frac{1}{2m}\sum_v\deg(v)\log(\deg(v)-1)
       =\frac{3(h-1)}{3h+k-3}>0.
\]

Put $t=\lfloor(g-1)/2\rfloor$.
The random walk of length $t$ has entropy
\[
 H_t=\log(2m)+(t-1)\sigma.
\]
Every such walk is a simple path, since $t<g$.
Two distinct such paths with the same ordered endpoints would give a cycle of length at most $2t<g$.
There are therefore at most $n^2$ possible walks.
Entropy is at most the logarithm of the support size, so
\[
 (t-1)\sigma\le2\log n-\log(2m)\le\log n.
\]
The final inequality uses $2m\ge n$.
Since $t\ge(g-2)/2$, this gives $g\le4+2\log n/\sigma$.
Substitute the formula for $\sigma$.
\end{proof}

\begin{corollary}[Small or polynomial-rank shores]
\label{cor:four-high-girth-shore-dichotomy}
Fix $c,C>0$.
For all sufficiently large $r$, the following holds.
Suppose a normalized $k$-pole has rank $h\ge1$, with
\[
 k\le C\log r,
 \qquad \operatorname{girth}(N)\ge c\log r.
\]
Then either
\[
 h\le1+\frac{4C}{3c}\log(3C\log r),
\]
or
\[
 h\ge r^{c/4}.
\]
The rank threshold depends only on $c,C$.
\end{corollary}

\begin{proof}
The case $h=1$ satisfies the first bound.
Assume $h\ge2$, and write $x=\log r$.
First suppose $h<k$.
Then $n\le3Cx$ and Lemma~\ref{lem:four-entropy-moore} gives
\[
 cx\le4+2\log(3Cx)
       +\frac{2Cx}{3(h-1)}\log(3Cx).
\]
For large $r$, the first two terms sum to at most $cx/2$.
Rearrangement gives the first alternative.

Now suppose $h\ge k$.
Then $n\le3h$ and $k/[3(h-1)]\le2/3$.
Consequently
\[
 cx\le4+\frac{10}{3}\log(3h),
 \qquad
 h\ge\frac{2^{-6/5}}{3}\,r^{3c/10}.
\]
For sufficiently large $r$, the last bound exceeds $r^{c/4}$.
\end{proof}

For a minimum cyclic cut in a cubic graph, every shore is connected and its ports are distinct.
Thus each shore is a normalized pole.
The girth of either shore is at least the ambient girth.
The corollary therefore applies directly when the cut order is $O(\log r)$.

For minimum-cut shores, an elementary argument gives a stronger direct conclusion.
It avoids any external theorem about atoms or cyclic parts.

\begin{corollary}[Every noncycle shore has exponential rank in its girth]
\label{cor:four-minimum-shore-exponential-rank}
Let $N$ be a normalized $k$-pole of rank $h\ge2$ and girth $g\ge\max\{k,3\}$.
Then
\[
 k\le6(h-1),\qquad
 h\ge1+2^{(g-22)/6}.
\]
Consequently, a shore of a minimum cyclic cut in a cubic graph is either a cycle or has rank exponential in the ambient girth.
\end{corollary}

\begin{proof}
Every graph of rank at least two contains two distinct cycles.
If they are edge-disjoint, their union has at least $2g$ edges.
Otherwise their union contains a theta graph.
Its three cycles have total length twice its number of edges, so that number is at least $3g/2$.
In either case, writing $m=3h+k-3$,
\[
 g\le\frac{2m}{3}=2h+\frac{2k}{3}-2.
\]
Since $g\ge k$, it follows that $k\le6(h-1)$.
Therefore
\[
 n=2h+k-2\le8(h-1),\qquad
 \sigma=\frac{3(h-1)}{3(h-1)+k}\ge\frac13.
\]
Lemma~\ref{lem:four-entropy-moore} gives
\[
 g\le4+6\log n\le22+6\log(h-1).
\]
Rearranging proves the exponential bound.

A minimum cyclic-cut shore is a normalized pole, and its girth is at least the ambient girth.
The cut order is at most the ambient girth.
Thus the bound applies whenever its rank is at least two.
A connected normalized pole of rank one has only degree-two vertices, so it is a cycle.
\end{proof}

This is a structural reduction, not a cycle-count bound by itself.
The crossing-count application also needs the growing-terminal forest estimate and its uniform error.

\begin{corollary}[A stronger rank bound and a large port-free ball]
\label{cor:four-near-moore-shore-rank}
Under the hypotheses of Corollary~\ref{cor:four-minimum-shore-exponential-rank}, if $g\ge120$, then
\[
 h\ge2^{g/2-5}.
\]
Moreover, $N$ has a port-free radius-$a$ ball, where
\[
 a=\left\lfloor\frac g2-\log(192g)\right\rfloor.
\]
This ball contains at least $2^{g/2}/(192g)$ vertices.
\end{corollary}

\begin{proof}
The preceding corollary gives $h-1\ge2^{(g-22)/6}\ge g^2$ for $g\ge120$.
The final comparison holds at $120$, and its left-to-right ratio increases thereafter.
If $h\ge2^{g/2}$, the first assertion already holds.
Otherwise
\[
 n=2h+k-2\le2^{g/2+2},\qquad
 \frac{2k\log n}{3(h-1)}
 \le\frac{2g(g/2+2)}{3g^2}<1.
\]
Lemma~\ref{lem:four-entropy-moore} now gives $g<5+2\log n$.
Thus $n>2^{(g-5)/2}$.
For $g\ge120$, one has $g\le2^{(g-5)/2}/2$.
Since $k\le g$,
\[
 h=\frac{n-k+2}{2}>\frac14\,2^{(g-5)/2}>2^{g/2-5}.
\]

Any radius-$a$ ball in a graph of maximum degree three contains at most $3\cdot2^a-2$ vertices.
The union of those balls around all ports therefore has at most
\[
 3k2^a\le3g\frac{2^{g/2}}{192g}=2^{g/2-6}
\]
vertices.
But $n\ge h\ge2^{g/2-5}$, so some vertex lies outside that union.
Its radius-$a$ ball contains no port.
Since $2a<g$, its breadth-first expansion has at least $3\cdot2^a-2\ge2^{a+1}$ vertices.
The definition of $a$ gives $2^{a+1}\ge2^{g/2}/(192g)$, as required.
\end{proof}

\section{Higher-terminal forest stability at high girth}
\label{sec:path-high-girth-terminal-forests}

The root agent proposed using a breadth-first tree marker for all ports.
The localization argument of Fable then extends from one terminal path to any even terminal set.
A growing-window argument bounds each distance layer, while an entropy Moore bound supplies a large cubic ball.
All logarithms have base two.

\begin{lemma}[A marker with few vertices in each layer]
\label{lem:path-bfs-terminal-marker}
Let $N$ be a normalized $k$-pole with $k\ge2$, and let $T$ be any even subset of its ports.
Choose a breadth-first tree rooted at one port.
Let $P$ be the union of its root-to-port paths, as a vertex set.
Then $P$ is connected, contains every port, and meets each distance layer in at most $k$ vertices.
Moreover, if $J$ is a uniform $T$-join and $\xi$ is a uniform even subgraph, then
\[
 \Pr(J\text{ is a forest})
 \le\Pr(\xi-P\text{ is a forest}).
\]
\end{lemma}

\begin{proof}
Each root-to-port path has at most one vertex in each distance layer.
Their union gives the asserted layer bound and contains all degree-two vertices.
The union of their tree edges is connected, so it contains a $T$-join $F_0$.
The map $\xi\mapsto\xi\mathbin\triangle F_0$ is a bijection onto the affine $T$-join space.
Its restriction outside $P$ agrees with that of $\xi$.
A forest image therefore implies the forest event outside $P$.
\end{proof}

\begin{lemma}[Layer width from a thin connected marker]
\label{lem:path-thin-marker-layer-width}
There is an absolute constant $K$ with the following property.
Let $N$ be connected and subcubic, and let $P$ be a connected marked set containing every vertex of degree below three.
Suppose $P$ meets each distance layer from a fixed root in at most $k$ vertices, where $k\ge1$.
If
\[
 \Pr(\xi_N-P\text{ is a forest})\ge\frac12+\varepsilon,
 \qquad 0<\varepsilon\le\frac12,
\]
then every full distance layer has size at most
\[
 K\left[
 \varepsilon^{-2}\bigl(\log(4/\varepsilon)\bigr)^{11}
 +(k+1)^2\right].
\]
\end{lemma}

\begin{proof}
Put $\delta=\varepsilon/2$, and use $R_\delta$ from
Lemma~\ref{lem:path-polynomial-local-threshold}.
Set
\[
 \Theta=\max\{2R_\delta,\,2^{100}(k+1)^2\}.
\]
We first exclude every set $R\subseteq V(N)\setminus P$ satisfying
\[
 |R|\ge\Theta,
 \qquad
 |\delta_N R|\le\frac{|R|}{4(\log|R|)^2}.
\]
Indeed, let $C_0$ be the component of $N-R$ containing $P$, and put
$R^+=V(N)\setminus V(C_0)$.
This set contains $R$, avoids $P$, and has connected complement.
Its boundary is contained in $\delta_N R$.
All its vertices have degree three in $N$.
Lemma~\ref{lem:path-disconnected-localization} therefore gives a localized graph of rank
\[
 \rho=\frac{|R^+|+|\delta_N R^+|}{2}+1\ge|R|/2\ge R_\delta.
\]
Its marked cycle has at most $\rho/(\log\rho)^2$ vertices, by exactly the monotonicity and padding comparison in Lemma~\ref{lem:path-direct-good-set}.
The forest event outside $P$ implies the forest event on $R^+$.
Localization and the definition of $R_\delta$ bound the latter probability by
$1/2+\delta$, a contradiction.

Write $m_j$ for the number of unmarked vertices in layer $j$, with $m_j=0$ outside the layer range.
For the union $R(i,i')$ of its unmarked layers from $i$ through $i'$, let $m$ be its order.
Edges join the same or consecutive layers.
The boundary satisfies
\[
 |\delta_N R(i,i')|
 \le3m_{i-1}+3m_{i'+1}+3k(i'-i+3).
\]
The last term bounds all incidences to marked vertices in layers $i-1$ through $i'+1$.

Suppose $m_j\ge\Theta$.
Start with this layer and include a neighboring layer whenever its mass exceeds
$m/(48(\log m)^2)$, where $m$ is the current window mass.
This process stops because the graph is finite.
Let $\mu$ be the final mass and put $L=\log\mu\ge100$.
Every step multiplies the current mass by more than $1+1/(48L^2)$.
The number of steps is at most
\[
 \frac{\ln\mu}{\ln(1+1/(48L^2))}
 \le96L^2\ln\mu<67L^3.
\]
Both remaining neighbor masses are at most $\mu/(48L^2)$.
Consequently
\[
 |\delta_N R(i,i')|
 \le\frac{\mu}{8L^2}+3k(67L^3+3)
 \le\frac{\mu}{4L^2}.
\]
For completeness, the last inequality follows from
\[
 24kL^2(67L^3+3)\le2000kL^5\le2^L.
\]
The first comparison holds for $L\ge100$.
For the second, $\mu\ge2^{100}(k+1)^2$ gives $k\le2^{L/2-50}$, and
$2000L^5\le2^{L/2+50}$ for $L\ge100$.
The latter inequality holds at $100$, and its right-to-left ratio increases thereafter.
Thus the stopped window is an excluded set, a contradiction.

Every layer therefore has fewer than $\Theta$ unmarked vertices and at most $k$ marked vertices.
Finally,
\[
 R_\delta\le K_0\delta^{-2}\bigl(\log(2/\delta)\bigr)^{11}.
\]
Absorb this bound and the fixed coefficient $2^{100}$ into an absolute $K$.
\end{proof}

\begin{theorem}[Uniform affine forest bound for high-girth poles]
\label{thm:path-high-girth-terminal-forests}
There is an absolute constant $C$ such that the following holds.
Let $N$ be a normalized $k$-pole of rank $h\ge2$ and girth $g$, with
\[
 2\le k\le g.
\]
For every even subset $T$ of its ports,
\[
 \frac{\#\{T\text{-joins in }N\text{ that are forests}\}}{2^h}
 \le\frac12+C g^{13/2}2^{-g/4}.
\]
In particular, the excess tends to zero as $g\to\infty$, uniformly in the rank, the port count, and the even terminal set.
\end{theorem}

\begin{proof}
It suffices to prove the assertion for all sufficiently large $g$.
An increase in $C$ then covers the remaining finitely many girths.
Corollary~\ref{cor:four-near-moore-shore-rank} gives
\[
 n=|V(N)|\ge2^{g/2-5}.
\]
Put
\[
 s=\left\lfloor\frac g2-\log k-10\right\rfloor.
\]
For large $g$, $s\ge1$ and $2s<g$.
The union of the radius-$s$ balls about the ports has at most
\[
 3k2^s\le3\,2^{g/2-10}<n
\]
vertices.
Choose a vertex $v$ outside that union.
The radius-$s$ ball about $v$ contains no port.
Its breadth-first expansion is cubic and has no repeated vertex before radius $s$, since $2s<g$.
It therefore contains at least
\[
 2^s\ge\frac{2^{g/2}}{2^{11}k}
\]
vertices.

Use the breadth-first root and marker from Lemma~\ref{lem:path-bfs-terminal-marker}.
Distances from this root vary by at most $2s$ across the ball about $v$.
Thus that ball meets at most $2s+1\le g$ root-distance layers.
Some layer has size at least
\begin{equation}\label{eq:path-high-girth-layer-lower}
 \frac{2^{g/2}}{2^{11}kg}
 \ge\frac{2^{g/2}}{2^{11}g^2}.
\end{equation}

Suppose the affine forest proportion is at least $1/2+\varepsilon$.
Translation and Lemma~\ref{lem:path-thin-marker-layer-width} bound every layer by
\[
 K\left[
 \varepsilon^{-2}\bigl(\log(4/\varepsilon)\bigr)^{11}
 +(k+1)^2\right].
\]
Set $\varepsilon=Cg^{13/2}2^{-g/4}$.
For sufficiently large $g$, this lies in $(0,1/2]$ and
$\log(4/\varepsilon)\le g$.
The last layer bound is then at most
\[
 \frac{K}{C^2}\frac{2^{g/2}}{g^2}+K(g+1)^2.
\]
Choose $C$ with $K/C^2\le2^{-13}$.
For sufficiently large $g$, the polynomial second term is also at most
$2^{-13}2^{g/2}/g^2$.
The resulting upper bound is strictly smaller than
\eqref{eq:path-high-girth-layer-lower}, a contradiction.
This proves the asserted uniform estimate.
\end{proof}

The hypothesis $h\ge2$ is essential for this argument.
A cycle with every vertex a port has rank one and girth $k$.
Every nonempty even-terminal affine class in that pole consists entirely of forests.

\begin{corollary}[A quarter-space bound across a nontrivial minimum cut]
\label{cor:path-high-girth-noncycle-minimum-cut}
There is an absolute constant $C_1$ such that the following holds for all sufficiently large girths $g$.
Let $G$ be a connected loopless cubic multigraph of rank $r$ and girth $g$.
Suppose a minimum cyclic cut has order $k$ and neither of its shores is a cycle.
Then
\[
 \frac{c(G)}{2^r}
 \le\frac14+C_1 g^{13/2}2^{-g/4}.
\]
Consequently, any sequence of such graphs with ranks tending to infinity, no cyclic cut of size at most four, and cycle counts above $2^{r-1}$ must eventually have the following property:
its minimum cyclic cut order equals its girth, and every minimum cyclic cut has a shortest cycle as one shore.
\end{corollary}

\begin{proof}
At sufficiently large girth, $G$ has a cyclic cut and its minimum order satisfies $k\le g$.
The shores of a minimum cyclic cut are connected, and its boundary edges have distinct endpoints on each shore.
Write $N_1,N_2$ for these normalized $k$-poles, and $h_1,h_2$ for their ranks.
A rank-one normalized pole is a cycle.
Thus $h_1,h_2\ge2$, and
\[
 r=h_1+h_2+k-1.
\]
The near-Moore rank bound gives $h_i\ge2^{g/2-5}$.
If $k=1$, no cycle crosses the cut; the internal-cycle estimate below proves the assertion directly.
Assume henceforth that $k\ge2$.

A cycle crossing the cut uses a nonempty even subset $S$ of its edges.
Its restriction to each shore is an $S$-join forest.
Some pairs of such forests fail to form one cycle, so their product is an upper bound for the crossing count in this sector.
Apply Theorem~\ref{thm:path-high-girth-terminal-forests} to each shore, with the ambient girth lower bound $g$.
The proof of that theorem applies unchanged with a girth lower bound in place of the exact girth.
Writing $\eta_g=Cg^{13/2}2^{-g/4}$, each sector contributes at most
\[
 (1/2+\eta_g)^2 2^{h_1+h_2}.
\]
There are $2^{k-1}-1$ nonempty even subsets.
After normalization by $2^r$, all crossing cycles therefore contribute at most
$(1/2+\eta_g)^2$.

The cycles internal to the shores contribute at most
\[
 \frac{2^{h_1}+2^{h_2}}{2^r}
 =2^{-h_2-k+1}+2^{-h_1-k+1}.
\]
The rank lower bounds make this $o(\eta_g)$.
Absorb this term and $\eta_g^2$ into an absolute multiple of $g^{13/2}2^{-g/4}$.
This proves the first assertion.

For the final assertion, Corollary~\ref{cor:path-unbounded-girth-reduction} forces the girth to tend to infinity.
The first assertion then excludes any minimum cyclic cut with two noncycle shores.
Every remaining minimum cut has a cycle shore.
Its boundary order equals the length of that cycle, so it is at least the girth.
Since the minimum cyclic cut order is at most the girth, equality holds, and that shore is a shortest cycle.
\end{proof}

\section{A direct cycle deficit in a growing girth window}
\label{sec:random-girth-window}

Moore-ball packing makes the shortest-cycle prefix argument quantitative.
The estimate below does not require a restriction on cyclic cuts.
All logarithms are binary.

\begin{lemma}[A direct bound for the large prefix branches]
\label{lem:random-girth-prefix-path}
There are absolute constants $C,R$ with the following property.
Let $r\ge R$, put $x=\log r$, $z=\log x$, and set
\[
 g_0=\left\lceil4z+2\log z+12\right\rceil.
\]
Let $N$ be a normalized $k$-pole with rank $b\in[r/2,r]$, girth at least $g_0$, and $k\le3x$.
For any distinct ports $u,v$,
\[
 P_{uv}(N)\le
 \left(\frac12+C\frac{x^{11/2}}{\sqrt r}\right)2^b.
\]
\end{lemma}

\begin{proof}
The order of $N$ is $n=2b+k-2$.
Apply Lemma~\ref{lem:four-pole-moore-geodesic} with the integer girth lower bound $g_0$.
The terminal distance $\ell$ satisfies
\[
 \ell+1<g_0\left(k+1+n2^{-g_0/2}\right).
\]
For sufficiently large $r$,
\[
 g_0\le8z,\qquad
 2^{-g_0/2}\le\frac1{64x^2z},\qquad
 n\le3b.
\]
Hence
\[
 \ell+2\le8z(3x+1)+\frac{3b}{8x^2}+1
 \le\frac{b}{2x^2}.
\]
The last comparison is uniform for $b\ge r/2$.

Suppress every degree-two vertex other than $u,v$.
This preserves rank and terminal-path count, and it can only shorten terminal distance.
Call the resulting two-pole $N'$.
Parallel edges or loops do not affect these assertions or the multigraph forest theorem.
Add a new two-edge $u$--$v$ path to $N'$, obtaining a subcubic graph $H$ of rank $\rho=b+1$.
Let $A$ be a shortest terminal path in $N'$ and let $W$ be its union with the new path.
Then $W$ is a cycle and
\[
 |V(W)|\le\ell+2\le\frac{b}{2x^2}
 \le\frac{\rho}{(\log\rho)^2}.
\]
For the final inequality, use $\rho\le r+1$ and take $r$ sufficiently large.
Every vertex outside $W$ is cubic.

Each terminal path $Q$ of $N'$ gives two even subgraphs of $H$:
its union with the new path, and $Q\mathbin{\triangle}A$.
These $2P_{uv}(N')$ subgraphs are distinct.
Outside $V(W)$, each is a subgraph of $Q$ and is therefore a forest.
Consequently
\[
 \frac{P_{uv}(N)}{2^b}
 \le
 \Pr\bigl(\xi_H-V(W)\text{ is a forest}\bigr).
\]
The endpoint forest theorem, Theorem~\ref{thm:parallel-links-log-five-half-forest}, bounds the right side by
\[
 \frac12+O\!\left(\frac{(\log\rho)^{11/2}}{\sqrt\rho}\right)
 \le\frac12+C\frac{x^{11/2}}{\sqrt r}.
\]
All constants are absolute.
\end{proof}

\begin{theorem}[The quantitative girth window]
\label{thm:random-girth-window}
There are absolute constants $C,R$ such that every connected loopless cubic multigraph $G$ of rank $r\ge R$ and girth $g$ satisfying
\[
 \left\lceil4\log\log r+2\log\log\log r+12\right\rceil
 \le g
\]
obeys
\[
 c(G)\le
 2^{r-1}-2^{r-g-1}
 +C\,2^r\frac{(\log r)^{11/2}}{\sqrt r}.
\]
No hypothesis on cyclic connectivity is required.
\end{theorem}

\begin{proof}
Choose a shortest cycle $C=v_0\cdots v_{g-1}v_0$.
The cubic Moore bound gives $g\le2\log(2r-2)+2\le3\log r$ for sufficiently large $r$.
Delete its edge $e=v_0v_{g-1}$ and write $b=r-1$.
Partition the remaining terminal paths by first divergence from the fixed path $C-e$.
Lemma~\ref{lem:path-shortest-cycle-prefix} gives, for each $1\le p\le g-1$, either no branch or a normalized pole of rank
\[
 b_p\le b-p.
\]
Its two terminals are distinct ports.
It has at most $g$ ports and girth at least $g$.
The last branch includes the pendant-endpoint deletion stated in that lemma.

Put $\varepsilon_r=C_1(\log r)^{11/2}/\sqrt r$, with $C_1$ from Lemma~\ref{lem:random-girth-prefix-path}.
If $b_p\ge r/2$, that lemma bounds the branch count by
\[
 (1/2+\varepsilon_r)2^{b_p}
 \le(1/2+\varepsilon_r)2^{b-p}.
\]
If $b_p<r/2$, the parity-class bound gives at most $2^{r/2}$ paths.
Including the fixed path $C-e$ and summing the geometric series therefore gives
\begin{align*}
 P_{v_0v_{g-1}}(G-e)
 &\le1+(1/2+\varepsilon_r)\sum_{p=1}^{g-1}2^{b-p}
       +(g-1)2^{r/2}\\
 &=1+(1/2+\varepsilon_r)(2^b-2^{b-g+1})
       +(g-1)2^{r/2}.
\end{align*}

The unrestricted endpoint cycle bound applies to the connected graph $G-e$:
\[
 c(G-e)\le2^{b-1}
 +O\!\left(2^b\frac{(\log b)^{11/2}}{\sqrt b}\right).
\]
Add this estimate to the terminal-path count.
The two half terms give $2^{r-1}-2^{r-g-1}$.
Since $g\le3\log r$, the term $1+(g-1)2^{r/2}$ is absorbed by the displayed endpoint error.
Increasing the absolute constant proves the theorem.
\end{proof}

\begin{corollary}[A strict deficit up to half the rank logarithm]
\label{cor:random-girth-window-strict}
There is an absolute constant $C_0$ such that, for sufficiently large $r$, every connected loopless cubic multigraph of rank $r$ and girth in the interval
\[
 \left\lceil4\log\log r+2\log\log\log r+12\right\rceil
 \le g
 \le\frac12\log r-\frac{11}{2}\log\log r-C_0
\]
satisfies
\[
 c(G)\le2^{r-1}-2^{r-g-2}<2^{r-1}.
\]
More generally, within the theorem's girth range, if
\[
 2^g\frac{(\log r)^{11/2}}{\sqrt r}\longrightarrow0,
\]
then the negative term is $(1-o(1))2^{r-g-1}$.
\end{corollary}

\begin{proof}
Choose $C_0$ so that $C2^{-C_0}\le1/4$, where $C$ is the theorem's error constant.
The upper girth bound then gives
\[
 C\frac{(\log r)^{11/2}}{\sqrt r}
 \le2^{-g-2}.
\]
Subtract this error from $2^{-g-1}$ after normalizing by $2^r$.
The limiting assertion follows from the same comparison.
\end{proof}

The remaining unrestricted problem includes girths between the available double-logarithmic ranges and girths above this window.
This theorem supplies a quantitative exclusion range; it does not prove the unrestricted sharp-order error.

\section{A larger marked set and a lower girth threshold}
\label{sec:random-larger-marker}

The endpoint forest proof permits a fixed multiple of $r/\log r$ marked vertices.
Only its initial spectral loss changes.
All logarithms below have base two.

\begin{theorem}[The endpoint estimate with a larger marked set]
\label{thm:random-larger-marker}
Put
\[
 A=6e^{168},\qquad a=\frac1{1000A},\qquad
 B=\lceil\log(10^6A)\rceil,\qquad M=64A^5 2^{B/2}.
\]
Let $\Gamma$ be a connected subcubic multigraph of rank $r$.
Let $P$ contain all vertices of degree below three.
Write $x=\log r$ and $z=\log x$, and assume
\[
 |P|\le a\,\frac r x,\qquad
 \operatorname{comp}(\Gamma[P])\le\frac x{z^2}.
\]
Then, for all sufficiently large $r$, a uniform even subgraph satisfies
\[
 \Pr(\xi_\Gamma-P\text{ is a forest})
 \le\frac12+M\frac{x^{11/2}}{\sqrt r}.
\]
The threshold is absolute.
\end{theorem}

\begin{proof}
If $P$ is empty, the probability is $2^{-r}$.
Otherwise use the parameters from Theorem~\ref{thm:parallel-links-log-five-half-forest}:
\[
 L=\lfloor Ax\rfloor,\quad
 t=\lfloor z-2\log z\rfloor,\quad
 D_0=\left\lceil\frac{8Ax}{z}\right\rceil,\quad
 D_1=\lfloor x-3z-B\rfloor.
\]
Assume for contradiction that the probability exceeds
$1/2+\varepsilon$, where $\varepsilon=Mx^{11/2}/\sqrt r$.
The nested short-cycle packings give
\[
 |S_0|=r^{o(1)},\qquad
 |S_1|\le12\,2^{-B}r/x.
\]
These estimates use the forest excess and the boundary sizes of cycles outside $P$.
They do not use $|P|\le r/x^2$.

The degree sum gives $|V(\Gamma)|=2r+O(|P|+1)$.
Thus $v_1=|V(\Gamma-(P\cup S_1))|=2r+o(r)$, and in particular $v_1\ge r$ for large $r$.
Its degree deficit satisfies $d_1\le3(|P|+|S_1|)$.
The only changed spectral estimate is
\[
 \frac{Ld_1}{2v_1}
 \le\frac{3Aa}{2}+18A2^{-B}
 \le\frac3{2000}+\frac{18}{10^6}
 <\frac1{100}.
\]
Consequently the same even trace sum supplies $\mu_1>20$.

The link estimate uses the number of connected components of $P$, not its order.
It therefore still gives at most $3Ax/z$ links at either girth scale.
The exceptional-pair error remains $O(L^3 2^{-Ax/z})=o(1)$.
The two-link moment bounds, the anchor-component argument, and weighted maximal disjointness consequently give
\[
 T\ge\frac{2^{1-B/2}}{A^3}\frac{\sqrt r}{x^{9/2}}
\]
nested separators.

Use full blocks with $\lceil16A^2x\rceil+2$ gaps, as in the endpoint proof.
Their orders are at least $8L^2$, and their boundaries have size at most $2L$.
The number discarded is at most
\[
 L+|S_0|+\operatorname{comp}(\Gamma[P])=r^{o(1)}.
\]
In particular, enlarging a connected component of $P$ cannot create extra discarded blocks:
that component cannot cross a block boundary, whose vertices lie on separator cycles disjoint from $P$.
Every retained block has degree deficit at most $2L$, girth greater than $D_0$, and order at most $3r$.
The one-block spectral estimate therefore gives
\[
 \mu_0\ge c_B\frac{\sqrt r}{x^{11/2}},
 \qquad c_B=\frac{2^{-B/2}}{64A^5}.
\]
The inverse moment estimate gives instead
\[
 \mu_0\le\frac{1+o(1)}{4\varepsilon}
 =(1+o(1))\frac{c_B}{4}\frac{\sqrt r}{x^{11/2}},
\]
a contradiction.
This lists every place where the marked-set hypothesis enters the proof.
\end{proof}

\begin{corollary}[Improved girth window]
\label{cor:random-improved-girth-window}
There are absolute constants $C_*,C,C_0,R$.
Let $G$ be a connected loopless cubic multigraph of rank $r\ge R$.
If its girth satisfies
\[
 g\ge
 \left\lceil
 2\log\log r+2\log\log\log r+C_*
 \right\rceil
\]
satisfies
\[
 c(G)\le
 2^{r-1}-2^{r-g-1}
 +C\,2^r\frac{(\log r)^{11/2}}{\sqrt r}.
\]
If also
\[
 g\le\frac12\log r-\frac{11}{2}\log\log r-C_0,
\]
then $c(G)\le2^{r-1}-2^{r-g-2}$.
Neither assertion requires a cyclic-connectivity assumption.
\end{corollary}

\begin{proof}
Use the proof of Theorem~\ref{thm:random-girth-window}, replacing its large-branch lemma as follows.
Put $x=\log r$, $z=\log x$, and take
\[
 C_*=\left\lceil2\log(1000/a)\right\rceil,\qquad
 g_*=\lceil2z+2\log z+C_*\rceil.
\]
For large $r$, $g_*\le6z$ and
\[
 2^{-g_*/2}\le\frac{a}{1000xz}.
\]
Let a branch pole have rank $b\in[r/2,r]$ and at most $3x$ ports.
Its order is at most $3b$.
Moore-ball packing bounds its terminal distance by
\[
 \ell+2
 \le6z(3x+1)+\frac{18ab}{1000x}+1
 \le\frac{ab}{2x}
\]
for all sufficiently large $r$, uniformly over these branches.
Suppress its other ports, and add the new two-edge terminal path as before.
The resulting rank is $\rho=b+1$.
Its marked cycle is connected, contains every vertex of degree below three, and has order at most
\[
 \frac{ab}{2x}\le a\frac{\rho}{\log\rho}.
\]
Theorem~\ref{thm:random-larger-marker} gives the same branch error
$O(x^{11/2}/\sqrt r)$.
The low-rank branch bound, geometric prefix sum, and unrestricted cycle error are unchanged.
They prove the finite estimate.
Choose $C_0$ as in Corollary~\ref{cor:random-girth-window-strict} to obtain the strict deficit.
\end{proof}

\section{Faster cut induction when girth exceeds the cut order}
\label{sec:path-girth-surplus-cuts}

The cut induction can use the endpoint forest error directly when its selected edges have short alternate paths.
A girth surplus guarantees such a choice at every deletion step.
All logarithms have base two.

\begin{lemma}[Moore balls along a terminal geodesic]
\label{lem:path-geodesic-moore-balls}
Let $N$ be a connected graph on $n$ vertices with distinct degree-two terminals $u,v$ and every other vertex cubic.
Suppose its girth is at least an integer $h\ge5$.
Then
\[
 \operatorname{dist}_N(u,v)
 \le2hn2^{-h/2}+h+1.
\]
\end{lemma}

\begin{proof}
Put $a=\lfloor(h-1)/2\rfloor$, and choose a shortest terminal path of length $\ell$.
If $\ell\le2a+1$, the asserted bound is immediate.
Otherwise choose centers along this path, starting at distance $a+1$ from $u$, spaced $2a+1$ apart, and stopping before distance $a$ from $v$.
The radius-$a$ balls around these centers are disjoint, because every subpath of the geodesic is shortest.
They avoid both degree-two terminals.

A breadth-first tree in any such ball has at least $1+3(2^a-1)\ge2^a$ vertices.
A repeated vertex in that tree would give a cycle of length at most $2a<h$.
The number $t$ of chosen centers therefore satisfies $t\le n2^{-a}$.
The choice of centers also gives
\[
 \ell\le(2a+1)t+2a+2
 \le hn2^{-a}+h+1.
\]
Finally $a\ge(h-2)/2$, so $2^{-a}\le2\,2^{-h/2}$.
\end{proof}

\begin{lemma}[A short alternate path gives the endpoint error]
\label{lem:path-short-alternate-endpoint}
Let $G$ be a connected cubic multigraph of rank $s$, and let $e=uv$ be a nonbridge edge.
If
\[
 \operatorname{dist}_{G-e}(u,v)+2\le\frac{s}{(\log s)^2},
\]
then, for all sufficiently large $s$,
\[
 P_{uv}(G-e)
 \le\left(\frac12+M\frac{(\log s)^{11/2}}{\sqrt s}\right)2^{s-1},
\]
where $M$ is an absolute constant.
\end{lemma}

\begin{proof}
Let $A$ be a shortest terminal path in $N=G-e$.
Add a new two-edge path $u-z-v$, obtaining a graph $H$ of rank $s$.
Mark $P=V(A)\cup\{z\}$.
This set is connected and contains every vertex of degree below three.
Its size satisfies the endpoint forest hypothesis.

For each simple terminal path $Q$ in $N$, the edge sets
\[
 Q\cup\{uz,zv\},\qquad Q\mathbin\triangle A
\]
are even subgraphs of $H$ whose restrictions outside $P$ are forests.
They give $2P_{uv}(N)$ distinct states, separated by whether the new two-edge path is present.
The endpoint forest theorem therefore gives
\[
 \frac{P_{uv}(N)}{2^{s-1}}
 \le\Pr(\xi_H\text{ is a forest outside }P)
 \le\frac12+M\frac{(\log s)^{11/2}}{\sqrt s}.
\]
\end{proof}

\begin{lemma}[Selecting fast edges through the cut reduction]
\label{lem:path-girth-surplus-edge-selection}
Fix $\eta>0$.
Let $G$ have rank $r$, girth $g$, and minimum cyclic cut order $k\ge5$.
Suppose
\[
 g-k\ge(4+\eta)\log\log r.
\]
For all sufficiently large $r$, one can carry out the $k-5$ cut-deletion steps of Proposition~\ref{prop:path-finite-cut-recurrence} so that every selected edge satisfies Lemma~\ref{lem:path-short-alternate-endpoint}.
The rank threshold may depend on $\eta$.
\end{lemma}

\begin{proof}
Retain one original minimum cut and its two shores throughout the reduction.
After $i$ steps, its remaining edges form a minimum cut of order $k-i$.
Each shore has lost exactly $i$ original port vertices by suppression.
Every cycle inside a current shore lifts to an original shore cycle by restoring at most $i$ vertices.
Thus every such cycle has length at least $g-i$.

Suppose the current graph has a cycle of length less than $g-i$.
It must cross the retained cut.
Choose an edge of that cut on the cycle.
Its alternate endpoint distance is at most $g-i-1$.
The original Moore bound gives $g=O(\log r)$, so this distance meets the hypothesis of Lemma~\ref{lem:path-short-alternate-endpoint}.

Otherwise the current graph has girth at least $g-i$.
This is at least $(4+\eta)\log\log r$.
Choose any edge of its retained cut and delete it.
The resulting connected graph has only its two terminals of degree two.
Apply Lemma~\ref{lem:path-geodesic-moore-balls} with
\[
 h=\lfloor(4+\eta)\log\log r\rfloor.
\]
It bounds the alternate distance by
\[
 O_\eta\!\left(
 \frac{r\log\log r}{(\log r)^{2+\eta/2}}
 +\log\log r\right)
 =o\!\left(\frac r{(\log r)^2}\right).
\]
This also meets the short-alternate-path hypothesis.

The original Moore bound also gives $k\le g=O(\log r)$ at these ranks.
All current ranks are therefore $r-O(\log r)$.
Both estimates above hold uniformly over the deletion steps.
\end{proof}

\begin{theorem}[A growing cut order with a girth surplus]
\label{thm:path-girth-surplus-growing-cuts}
Fix $\eta>0$.
Let connected loopless cubic multigraphs have ranks $r\to\infty$, girths $g$, and minimum cyclic cut orders $k\ge5$.
Assume
\[
 g-k\ge(4+\eta)\log\log r,
 \qquad
 \frac{2^k(\log r)^{11/2}}{\sqrt r}\longrightarrow0.
\]
Then
\[
 c(G)\le2^{r-1}-(1-o(1))2^{r-k-1}<2^{r-1}.
\]
Moreover, there is a constant $C_\eta$ such that, for sufficiently large $r$, the first condition together with
\[
 k\le\frac12\log r-\frac{11}{2}\log\log r-C_\eta
\]
implies
\[
 c(G)\le2^{r-1}-2^{r-k-2}.
\]
\end{theorem}

\begin{proof}
Use the edge choices from Lemma~\ref{lem:path-girth-surplus-edge-selection}.
Repeat the exact finite cut recurrence, replacing its worst-case path excesses by the selected-edge estimates.
Their weighted sum is at most
\[
 (M+o(1))\frac{(\log r)^{11/2}}{\sqrt r},
\]
because the geometric weights sum to less than one and all current ranks are $r-O(\log r)$.
The base error remains $2^{-(k-5)}a_5(r-k+5)$.
Comparing these two errors with $2^{-k-1}$ proves the first assertion.

For the endpoint assertion, choose $C_\eta$ large enough that the selected-edge error is at most $2^{-k-3}$.
At sufficiently large rank, $a_5(r-k+5)\le1/256$ uniformly in this range.
Its weighted contribution is also at most $2^{-k-3}$.
The remaining deficit is therefore at least $2^{-k-2}$ after normalization by $2^r$.
\end{proof}

This result does not assume that arbitrary minimum-cut edges have short alternate paths.
It uses either a short crossing cycle or the Moore-ball bound at each step.
It leaves open the regime where the girth is close to the minimum cyclic cut order.

\begin{corollary}[A girth-dependent error throughout the cut induction]
\label{cor:path-girth-weighted-cut-error}
There is an absolute constant $C$ such that the following holds for all sufficiently large ranks $r$.
Let $G$ be a connected loopless cubic multigraph of rank $r$, girth $g$, and minimum cyclic cut order $k\ge5$.
Then
\begin{align*}
 \frac{c(G)}{2^r}
 \le{}&\frac12-2^{-k-1}
 +2^{-(k-5)}a_5(r-k+5)\\
 &+C\left(g^{11/2}2^{-g/4}
       +\frac{(\log r)^{11/2}}{\sqrt r}\right).
\end{align*}
In particular, a sequence of such graphs satisfies
\[
 c(G)\le2^{r-1}-(1-o(1))2^{r-k-1}<2^{r-1}
\]
if
\[
 2^k g^{11/2}2^{-g/4}\longrightarrow0,
 \qquad
 \frac{2^k(\log r)^{11/2}}{\sqrt r}\longrightarrow0.
\]
Here $a_5$ is the base error in Proposition~\ref{prop:path-finite-cut-recurrence}.
\end{corollary}

\begin{proof}
As before, retain an original minimum cut during the smoothing process.
The Moore bound and the cyclic cut around a shortest cycle give
$5\le k\le g=O(\log r)$ for all sufficiently large $r$.
At stage $i$, every cycle internal to a retained shore has length at least $g-i$.

If a shorter cycle exists, choose an edge of the retained cut on that cycle.
Its alternate path has length $O(\log r)$.
Lemma~\ref{lem:path-short-alternate-endpoint} bounds its normalized path excess by
\[
 O\!\left(\frac{(\log r)^{11/2}}{\sqrt r}\right),
\]
uniformly over these stages.
Otherwise the current graph, and hence its edge-deleted two-pole, has girth at least $g-i$.
Proposition~\ref{prop:four-girth-path-excess} bounds the selected-edge path excess by
\[
 C_0(g-i)^{11/2}2^{-(g-i)/4}.
\]
All stages in the recurrence have $0\le i\le k-6$, so $g-i\ge6$.
The bound therefore applies at every such stage.

Insert these choices into the exact recurrence.
The geometric weights on the short-path errors sum to less than one.
The other errors have total at most
\[
 \sum_{i=0}^{k-6}2^{-i-1}C_0(g-i)^{11/2}2^{-(g-i)/4}
 \le\frac{C_0}{2}g^{11/2}2^{-g/4}
       \sum_{i\ge0}2^{-3i/4}.
\]
Absorb the convergent series and the short-path constant into $C$.
This proves the finite bound.
For the final assertion, $r-k+5\to\infty$, and the weighted base error divided by $2^{-k-1}$ is $64a_5(r-k+5)=o(1)$.
The two stated conditions make the remaining errors $o(2^{-k})$.
\end{proof}

\begin{corollary}[The improved girth surplus]
\label{cor:path-two-log-girth-surplus}
Theorem~\ref{thm:path-girth-surplus-growing-cuts} remains valid if its girth-surplus hypothesis is replaced by
\[
 g-k\ge(2+\eta)\log\log r,
\]
for any fixed $\eta>0$.
The constants and rank thresholds may depend on $\eta$.
\end{corollary}

\begin{proof}
Use the larger-marker endpoint theorem, Theorem~\ref{thm:random-larger-marker}, with its absolute constant $a>0$.
The proof of Lemma~\ref{lem:path-short-alternate-endpoint} now applies whenever the alternate distance plus two is at most
\[
 a\,\frac{s}{\log s}.
\]
The marked cycle is connected, so its component condition holds automatically.
The normalized path error remains $O((\log s)^{11/2}/\sqrt s)$.

Repeat the selected-edge argument with
$h=\lfloor(2+\eta)\log\log r\rfloor$.
A short crossing cycle still gives an alternate path of length $O(\log r)$.
In the other case, the Moore-ball estimate bounds that distance by
\[
 O_\eta\!\left(
 \frac{r\log\log r}{(\log r)^{1+\eta/2}}
 +\log\log r\right)
 =o(r/\log r).
\]
It therefore meets the larger-marker threshold for all sufficiently large $r$.
Every subsequent step of the recurrence and both error comparisons are unchanged.
\end{proof}

\section{Affine terminal forests at high girth}
\label{sec:random-terminal-forest-girth}

The larger-marker estimate also applies to a union of terminal paths.
This gives a quantitative bound for an entire affine terminal-join class.
All logarithms have base two.

\begin{lemma}[Translation by a terminal forest]
\label{lem:random-terminal-forest-translation}
Let $N$ be a normalized $k$-pole, where $k\ge2$ is even, and let $T$ be its set of ports.
Let $F_0$ be a $T$-join with the fewest edges, and put $P=V(F_0)$.
Then $F_0$ consists of $k/2$ vertex-disjoint paths.
For a uniform $T$-join $J$ and a uniform even subgraph $\xi$ of $N$,
\[
 \Pr(J\text{ is a forest})
 \le\Pr(\xi-P\text{ is a forest}).
\]
Moreover, if every terminal pair has distance at most $\ell$, then
\[
 |P|\le\frac k2(\ell+1),\qquad
 \operatorname{comp}(N[P])\le\frac k2.
\]
\end{lemma}

\begin{proof}
A cycle in $F_0$ could be deleted while preserving its boundary, contradicting minimality.
At a port the selected degree is one; at every other vertex it is zero or two.
Thus $F_0$ is a union of $k/2$ disjoint terminal paths.
In particular, $P$ contains all ports.

Pair the terminals arbitrarily and choose a shortest path for each pair.
Their symmetric difference is a $T$-join with at most $k\ell/2$ edges.
Minimality gives $|E(F_0)|\le k\ell/2$, and hence
$|P|=|E(F_0)|+k/2\le k(\ell+1)/2$.
The induced graph on $P$ may join some of these paths, so it has at most $k/2$ components.

The map $\xi\mapsto F_0\mathbin\triangle\xi$ is a bijection from the even space to the affine $T$-join space.
Outside $P$, its image has exactly the same edges as $\xi$.
If the image is a forest, its restriction outside $P$ is a forest.
This proves the probability inequality.
\end{proof}

\begin{theorem}[The affine forest bound at high girth]
\label{thm:random-terminal-forest-girth}
There are absolute constants $C_*,M,R$ such that the following holds.
Let $N$ be a normalized even-$k$ pole of rank $h\ge R$.
Put $x=\log h$ and $z=\log x$.
Assume
\[
 2\le k\le\frac{2x}{z^2},
 \qquad
 \operatorname{girth}(N)\ge
 \left\lceil2z+2\log z+2\log k+C_*\right\rceil.
\]
Then the proportion of forests in its affine all-port join space satisfies
\[
 \frac{\#\{T\text{-join forests in }N\}}{2^h}
 \le\frac12+M\frac{x^{11/2}}{\sqrt h}.
\]
No condition on bridges or cyclic connectivity is imposed.
\end{theorem}

\begin{proof}
Take $a=1/(1000A)$ from Theorem~\ref{thm:random-larger-marker}, and choose
$C_*=\lceil2\log(1000/a)\rceil$.
Let
\[
 g_*=\lceil2z+2\log z+2\log k+C_*\rceil.
\]
Uniformly over the allowed values of $k$, for sufficiently large $h$,
\[
 g_*\le8z,\qquad
 2^{-g_*/2}\le\frac{a}{1000kxz},\qquad
 |V(N)|=2h+k-2\le3h.
\]
Moore-ball packing bounds every terminal distance by a number $\ell$ with
\[
 \ell+1<8z\left(k+1+\frac{3ah}{1000kxz}\right).
\]
The preceding lemma therefore gives a marker satisfying
\[
 |P|\le4kz(k+1)+\frac{12ah}{1000x}
 \le a\frac h x
\]
for sufficiently large $h$.
The last comparison is uniform, because $k\le2x/z^2$ makes its first term polynomial in $x$, while $h=2^x$.
Also $\operatorname{comp}(N[P])\le k/2\le x/z^2$, and $P$ contains every vertex of degree below three.

Apply Theorem~\ref{thm:random-larger-marker} and the translation inequality.
Its constant $M$ and a sufficiently large absolute threshold work uniformly over $k$.
\end{proof}

For each fixed even $k$, the girth threshold is
$2\log\log h+2\log\log\log h+O_k(1)$.
In this range, a full affine terminal-forest space is impossible at large rank.
This does not establish a uniform non-saturation gap at fixed girth or a small-cut characterization of full four-terminal forest spaces.

\begin{theorem}[One marker for every terminal sector]
\label{thm:random-all-terminal-sectors-girth}
There are absolute constants $C_*,M,R$ with the following property.
Let $N$ be a normalized $k$-pole of rank $h\ge R$, put $x=\log h$, $z=\log x$, and assume
\[
 2\le k\le x^2,\qquad
 \operatorname{girth}(N)\ge
 \left\lceil2z+2\log z+2\log k+C_*\right\rceil.
\]
For every even subset $S$ of its ports, the number $F_S(N)$ of forests with odd vertex set $S$ satisfies
\[
 F_S(N)\le
 \left(\frac12+M\frac{x^{11/2}}{\sqrt h}\right)2^h.
\]
The estimate is uniform over all sectors, including when $k$ grows with $h$.
\end{theorem}

\begin{proof}
Take the same $a$ as in Theorem~\ref{thm:random-larger-marker} and
$C_*=\lceil2\log(1000/a)\rceil$.
The integer girth lower bound $g_*$ in the statement satisfies, uniformly for $k\le x^2$,
\[
 g_*\le8z,\qquad 2^{-g_*/2}\le\frac{a}{1000kxz}
\]
once $h$ is sufficiently large.
Also $|V(N)|\le3h$.
Consequently every terminal distance is at most a number $\ell$ satisfying
\[
 \ell+1<8z(k+1)+\frac{24ah}{1000kx}.
\]

Fix one port and join it to every other port by a shortest path.
Their union is connected and has at most $(k-1)\ell+1$ vertices.
Let $Q$ be a spanning tree of this union and put $P=V(Q)$.
Then $P$ contains all ports and is connected in $N$, while
\[
 |P|\le k(\ell+1)
 <8kz(k+1)+\frac{24ah}{1000x}
 \le a\frac h x.
\]
The final comparison is uniform for $k\le x^2$.

For any even port set $S$, the tree $Q$ has an $S$-join $J_S$.
Translation by $J_S$ maps the even space bijectively to the affine $S$-join space.
Outside $P$, the translated and original states agree.
Thus an $S$-join forest gives an even state that is a forest outside $P$.
Apply Theorem~\ref{thm:random-larger-marker}.
The same marker and the same probability bound work for every $S$.
\end{proof}

The same proof gives a more general sufficient condition, without the restriction $k\le(\log h)^2$.
If an integer girth lower bound $g_0\ge3$ satisfies
\[
 k g_0\left(k+1+(2h+k-2)2^{-g_0/2}\right)
 \le a\frac{h}{\log h},
\]
then the conclusion of Theorem~\ref{thm:random-all-terminal-sectors-girth} holds.
Indeed, Moore-ball packing bounds $k(\ell+1)$ by the left side, and the same connected geodesic marker applies.

\begin{lemma}[The growing-port aggregate]
\label{lem:random-growing-port-aggregate}
Let $Y$ be a normalized $k$-pole of rank $j$, where $k\ge2$, and set $\rho=j+k-1$.
If $\rho$ is sufficiently large and
\[
 2k-2\le a\frac{\rho}{\log\rho},
\]
where $a$ is the constant in Theorem~\ref{thm:random-larger-marker}, then
\[
 \sum_{S\text{ even}}F_S(Y)
 \le
 \left(\frac12+M\frac{(\log\rho)^{11/2}}{\sqrt\rho}\right)2^\rho.
\]
\end{lemma}

\begin{proof}
Take a disjoint tree with $k$ leaves and degree three at every nonleaf vertex; for $k=2$, take a single edge.
Add one edge from each leaf to a distinct port of $Y$.
Mark the new tree vertices.
The resulting subcubic graph has rank $\rho$.
The tree has $2k-2$ vertices, is connected, and contains every vertex of degree below three.
Restriction of uniform even states to $Y$ covers each even terminal-join class uniformly, with its $2^j$ states.
The event of a forest outside the marked tree therefore has probability
$2^{-\rho}\sum_S F_S(Y)$.
Apply the larger-marker theorem.
\end{proof}

\begin{corollary}[A quarter-space bound with a sufficiently large smaller shore]
\label{cor:random-high-girth-large-shores}
There are absolute constants $C_1,C_2$ with the following property.
Let connected loopless cubic graphs have ranks $r\to\infty$ and girths
\[
 g\ge
 \left\lceil4\log\log r+2\log\log\log r+C_1\right\rceil.
\]
Suppose a minimum cyclic cut has order $k$ and smaller shore rank
\[
 j\ge C_2 k\log(k+2).
\]
Then, uniformly over these graphs,
\[
 c(G)\le\left(\frac14+o(1)\right)2^r.
\]
In particular, these graphs are strictly below the half bound at sufficiently large rank.
The same conclusion holds with the smaller-shore hypothesis replaced by
$j\to\infty$ and $2k-2\le a(j+k-1)/\log(j+k-1)$.
\end{corollary}

\begin{proof}
Write $h\ge j$ for the shore ranks.
Minimality makes both shores connected.
It also makes their cut endpoints distinct: a vertex with two cut incidences is a leaf in its cyclic shore, and moving it to the other shore reduces the cut without destroying either cycle.
Thus the shores are normalized $k$-poles, with
\[
 r=h+j+k-1.
\]

At these ranks, a shortest cycle has cyclic complement.
Indeed, since $g\ge5$, it is induced and no outside vertex meets it twice.
Every component outside it consequently has minimum degree at least two.
Its boundary is a cyclic cut of size $g$.
The cubic Moore bound therefore gives
\[
 k\le g\le3\log r
\]
for large $r$.

Apply Theorem~\ref{thm:random-all-terminal-sectors-girth} to the shore of rank $h$.
Since $h\ge(r-k+1)/2$ and $k\le3\log r$, one has $h\ge r/3$ and $k\le(\log h)^2$ eventually.
Since $h\le r$ and $k\le3\log r$, the girth threshold of that theorem is at most
\[
 4\log\log r+2\log\log\log r+C_*+2\log3+1.
\]
Choose $C_1$ above this constant.
It follows that every even terminal sector in that shore has at most
$(1/2+o(1))2^h$ forests, uniformly.

Choose the absolute constant $C_2$ large enough that
$j\ge C_2 k\log(k+2)$ implies
$2k-2\le a(j+k-1)/\log(j+k-1)$ for all $k\ge2$.
For example, take $C_2\ge16$ with $C_2/(\log(2C_2)+4)\ge1/a$.
The function $s/\log s$ is increasing in the relevant range.
At $\rho_0=C_2 k\log(k+2)+k-1$ one has
\[
 \frac{\rho_0}{\log\rho_0}
 \ge
 \frac{C_2 k\log(k+2)}{\log(2C_2)+2\log(k+2)}
 \ge\frac{2C_2 k}{\log(2C_2)+4}
 \ge\frac{2k}{a},
\]
which proves the assertion.
Moreover $j\to\infty$ uniformly: if $k$ tends to infinity this follows from the rank hypothesis, while if $k$ stays bounded the shore order $2j+k-2$ must be at least the girth, which tends to infinity.
The growing-port aggregate consequently bounds the other shore's total forest mass by
$(1/2+o(1))2^{j+k-1}$.
For each fixed forest in that shore, at most the entire matching terminal sector in the first shore can combine with it to form one cycle.
Every cycle crossing the cut has a unique such pair of forest restrictions.
Hence the number of crossing cycles is at most
\[
 (1/2+o(1))2^h
 (1/2+o(1))2^{j+k-1}
 = (1/4+o(1))2^r.
\]
Cycles internal to the two shores number at most $2^h+2^j$.
Their ratio to $2^r$ tends to zero uniformly because $h\ge r/3$ and $j\to\infty$.
This proves the claim.
\end{proof}

\begin{proposition}[An exact criterion for full terminal-forest spaces]
\label{prop:random-full-terminal-forest-criterion}
Let $N$ be a normalized even-$k$ pole and $T$ its set of ports.
Every $T$-join is a forest if and only if every cycle $C$ of $N$ disjoint from $T$ leaves a component of $N-V(C)$ containing an odd number of terminals.
\end{proposition}

\begin{proof}
If every component of $N-V(C)$ has an even number of terminals, choose a terminal join in each component.
Their union with $C$ is a cyclic $T$-join.

Conversely, in a cyclic $T$-join choose one of its cycles $C$.
It avoids every port, whose selected degree is one.
Every vertex on $C$ has selected degree two, since the available degree is three and its parity is even.
No selected edge leaves $C$.
The remaining selected edges therefore form terminal joins in the components of $N-V(C)$.
Each such component has an even number of terminals.
\end{proof}

The standard-library checker \texttt{check\_random\_terminal\_forest\_girth.py}
verifies the translation, the common geodesic marker, and the saturation criterion
on 247 normalized poles.
It checks 8,312 affine states and 3,206 even terminal sectors.
These finite identities corroborate the implementation; the estimates above are proved independently of the checks.

\section{Fixed multiples of the girth and bounded-rank exceptions}
\label{sec:four-high-connectivity-weighted-kernel}

Suppressing the ports gives a cubic kernel with positive integer edge lengths.
A weighted Moore argument isolates the possible bounded-rank exceptions.
The path-stability agent proposed the walk argument below.
All logarithms have base two.

\begin{lemma}[A weighted-kernel Moore bound]
\label{lem:four-weighted-kernel-moore}
Let $N$ be a normalized $k$-pole of rank $h\ge2$ and girth $g\ge3$.
Then, with $t=\lceil\log(8h)\rceil$,
\[
 g\le4t\left(1+\frac{k}{3(h-1)}\right).
\]
\end{lemma}

\begin{proof}
Suppress all degree-two vertices.
Since $N$ is connected and has rank at least two, the resulting kernel is a connected cubic multigraph.
Loops and parallel edges are allowed.
It has $n_0=2h-2$ vertices and $m_0=3h-3$ edges.
Give each kernel edge the length of the path that it replaces.
These positive lengths sum to $m_0+k$.
The weighted girth of the kernel equals $g$.

There are $2m_0\,2^{t-1}$ directed non-backtracking walks of length $t$ in the kernel.
Choose one uniformly.
Every directed edge is uniform at every step, so its expected total weight is
\[
 t\left(1+\frac{k}{m_0}\right).
\]
At least half the walks have weight at most twice this expectation.
Their number is at least
\[
 m_0\,2^{t-1}\ge12h(h-1)>(2h-2)^2=n_0^2.
\]

If one such walk repeats a vertex, its edge support contains a cycle of no greater weight.
Otherwise these walks are simple paths.
Two must have the same ordered endpoints, so their union contains a cycle of weight at most the sum of their weights.
In either case, a cycle has weight at most $4t(1+k/m_0)$.
This proves the bound.
\end{proof}

\begin{proposition}[A bounded-rank or near-Moore alternative]
\label{prop:four-fixed-ratio-shore-rank}
Fix $C>0$, and put
\[
 H_C=\left\lceil64(C+1)\log(16(C+1))\right\rceil+1.
\]
If a normalized $k$-pole has rank $h\ge H_C$, girth at least $g$, and $k\le Cg$, then
\[
 h\ge2^{g/8-4}.
\]
For sufficiently large $g$, depending only on $C$, the stronger bound
\[
 h\ge2^{g/2-5}
\]
holds.
\end{proposition}

\begin{proof}
The choice of $H_C$ ensures
\[
 h-1>\frac{8C}{3}\,\lceil\log(8h)\rceil
 \qquad(h\ge H_C).
\]
Here is one explicit verification.
Put $D=C+1$ and $u=\log(16D)\ge4$.
At $h=H_C$, one has $h-1\ge64Du$ and $\log h+4\le4u$.
The function $(h-1)/(\log h+4)$ increases for $h\ge2$.
Also $\lceil\log(8h)\rceil\le\log h+4$.
These facts give the displayed inequality for all $h\ge H_C$.

Lemma~\ref{lem:four-weighted-kernel-moore} and $k\le Cg$ now give
\[
 g\le4t+\frac{4tCg}{3(h-1)}<4t+\frac g2.
\]
Thus $g<8t\le8\log h+32$, proving the first bound.

For sufficiently large $g$, that bound implies $h-1\ge(C+1)g^2$.
If $h\ge2^{g/2}$, the stronger assertion already holds.
Otherwise $n=2h+k-2\le2^{g/2+2}$ for large $g$, and
\[
 \frac{2k\log n}{3(h-1)}
 \le\frac{2Cg(g/2+2)}{3(C+1)g^2}<1.
\]
The entropy Moore bound therefore gives $g<5+2\log n$.
Since $n>2^{(g-5)/2}$ eventually exceeds $2Cg$, one has
\[
 h=\frac{n-k+2}{2}>\frac n4>2^{g/2-5}.
\]
\end{proof}

\begin{corollary}[Uniform terminal forests above the bounded-rank range]
\label{cor:four-fixed-ratio-terminal-forests}
For each fixed $C>0$, there is $B_C$ such that every normalized $k$-pole of rank $h\ge H_C$ with
\[
 2\le k\le Cg,
 \qquad \operatorname{girth}(N)\ge g
\]
satisfies, for every even port subset $T$,
\[
 \frac{F_T(N)}{2^h}
 \le\frac12+B_C g^{13/2}2^{-g/4}.
\]
\end{corollary}

\begin{proof}
Use Proposition~\ref{prop:four-fixed-ratio-shore-rank} in the proof of Theorem~\ref{thm:path-high-girth-terminal-forests}.
The port-free ball with radius $\lfloor g/2-\log k-10\rfloor$ still exists.
Its layer lower bound is at least $2^{g/2}/(2^{11}Cg^2)$.
The thin-marker layer upper bound has polynomial term at most $(Cg+1)^2$.
The same inversion gives the result after changing its constant by a factor depending only on $C$.
Increasing that constant handles the remaining finite range of $g$.
\end{proof}

\begin{corollary}[All fixed multiples of the girth]
\label{cor:four-fixed-ratio-cyclic-cuts}
Fix $C>0$.
Let connected cubic graphs have girths $g\to\infty$.
If a cut of order at most $Cg$ has two connected shores, each of rank at least $H_C$, then
\[
 c(G)\le\left(\frac14+o_C(1)\right)2^{\beta(G)}.
\]
\end{corollary}

\begin{proof}
Move vertices with two cut incidences across the cut until its endpoints are distinct on each shore.
Such a vertex is a leaf in its connected positive-rank shore; moving it preserves that rank, increases the other rank by one, and decreases the cut order by one.
A vertex with three cut incidences would be isolated in its shore, which is impossible.
The resulting connected normalized poles have ranks $h,j\ge H_C$ and cut order $k'\le Cg$.
Proposition~\ref{prop:four-fixed-ratio-shore-rank} makes both ranks at least $2^{g/2-5}$ for large $g$.
If $k'=1$, there are no crossing cycles and the conclusion follows from the internal-cycle bound.
Otherwise apply Corollary~\ref{cor:four-fixed-ratio-terminal-forests} to both shores and sum over the $2^{k'-1}-1$ nonempty even boundary subsets.
The crossing count is at most $(1/2+o_C(1))^2 2^{\beta(G)}$.
The internal counts, at most $2^h+2^j$, are negligible after normalization.
\end{proof}

Consequently, an unresolved above-half family with girth tending to infinity cannot have such a cut.
For every fixed $C$, each cut of order at most $Cg$ with connected shores must have a shore of rank below $H_C$.
This does not prove that an appropriate moderate cut exists in every remaining graph.

\section{The chain argument without short-cycle deletions}
\label{sec:random-no-deletion-chain}

At girth proportional to the logarithm of the rank, the graph itself supplies the girth needed by the parallel-link argument.
No packing estimate depending on the size of the forest excess is needed.
All logarithms have base two.

\begin{lemma}[Normalizing a cut between protected cyclic cores]
\label{lem:random-protected-core-cut}
Let $G$ be connected and cubic.
Let $A,B$ be disjoint connected subgraphs of minimum degree at least two.
Suppose a vertex partition separating $A$ from $B$ has at most $q$ crossing edges.
Then there is such a partition with at most $q$ crossing edges whose two shores are connected normalized poles.
Their ranks are at least $\beta(A)$ and $\beta(B)$.
\end{lemma}

\begin{proof}
Choose a minimum edge cut among partitions with all of $A$ on one side and all of $B$ on the other.
Each shore is connected.
Indeed, a component not containing its protected connected subgraph could be moved to the other shore, strictly reducing the cut.

A protected vertex has at least two incidences within its protected subgraph, so at most one cut incidence.
If an unprotected vertex had at least two cut incidences, moving that vertex to the opposite shore would reduce the cut.
Thus every vertex has at most one cut incidence.
The shores have degree two at their cut endpoints and degree three elsewhere.
Their ranks are at least the ranks of their respective protected subgraphs.
\end{proof}

\begin{lemma}[A positive forest excess gives two large shores]
\label{lem:random-no-deletion-chain}
Fix $c>0$.
For all sufficiently large $r$, the following holds.
Let $G$ be connected and cubic of rank $r$ and girth at least $c\log r$.
Let $P$ be a connected marked set with $|P|\le3\log r$.
If
\[
 \Pr(\xi_G-P\text{ is a forest})>\frac12,
\]
then $G$ has a cut of order at most $2A\log r$, where $A=6e^{168}$, with connected normalized shores of ranks at least
$r^{\min\{c,1\}/10}$.
The threshold depends on $c$, but not on the positive amount of forest excess.
\end{lemma}

\begin{proof}
Replace $c$ by $\min\{c,1\}$ if necessary.
Put $x=\log r$, $z=\log x$, and use
\[
 L=\lfloor Ax\rfloor,\qquad
 D=\lfloor cx/2\rfloor,\qquad
 t=\lfloor z-2\log z\rfloor,\qquad
 s=\lfloor(D-1)/2\rfloor.
\]
The graph $J=G-P$ has girth greater than $D$.
Thus the empty family is a maximal packing of cycles of length at most $D$ outside $P$.
In the notation of the endpoint proof, its deletion set is $S=\varnothing$.

The graph $J$ has $v=2r+O(x)$ vertices and degree deficit at most $3|P|=O(x)$.
Hence $Ld_J/(2v)=o(1)$.
The even spectral trace sum over $6x\le\ell\le L$, together with
$L^2 2^{-s}=o(1)$, gives
\[
 \mu=\sum_{\substack{C\subseteq J\\|C|\le L}}2^{-|C|}
 \ge\left(\frac14-o(1)\right)\ln(A/6)>20.
\]

The parallel-link estimate applies to this empty packing.
Every pair of disjoint cycles of length at most $L$ has at most
\[
 1+\frac{2L}{t+1}+2^{t+1}+578(L/D)^2
 \le\frac{3Ax}{z}
\]
links for sufficiently large $r$.
The exceptional-pair error in the established refined second-moment estimate of
\texttt{polynomial-cycle-error.pdf} is therefore
\[
 O\!\left(2^{3Ax/z}L^3 2^{-s}\right)
 =r^{-c/4+o(1)}=o(1).
\]
Only the first conclusion of that lemma is used.
It requires forest probability greater than one half and gives
$W_2>\mu^2-(1+o(1))\mu$.

Apply the anchor-component, far-pair, and weighted maximal-disjointness argument from the established endpoint forest proof in
\texttt{polynomial-cycle-error.pdf}.
For clarity, its output is a cycle $C\subseteq J$ of length at most $L$, a neighbor $u$ of $C$, and vertex-disjoint cycles
$C_1,\ldots,C_T$ in the component $H_0$ of $G-V(C)$ containing $P$.
These cycles separate $u$ from $P$ in $H_0$ and are nested in that order.
The weighted partner mass is at least $8/L$, so
\[
 T\ge\frac{8\,2^s}{L^3}
 \ge\frac{2}{A^3}\frac{r^{c/4}}{x^3}.
\]
The proof of this output uses no inverse bound on $\mu$ and no positive lower bound on the forest excess.

Choose a middle index $i$.
Let $U_i$ be the $u$-component of $H_0-V(C_i)$.
It contains $C_1,\ldots,C_{i-1}$.
The complementary induced subgraph $H_0-U_i$ is connected: it contains the connected cycle $C_i$, and every other component of $H_0-V(C_i)$ has an edge to that cycle.
It contains $C_i,\ldots,C_T$.
Thus the ranks of these two connected subgraphs are at least $i-1$ and $T-i+1$, respectively.

The boundary of $U_i$ in $G$ has at most $2L$ edges.
At most $L$ end on $C_i$, since a vertex of that cycle has at most one incidence outside it.
The remaining boundary edges end on the removed anchor cycle $C$, and there are at most $L$ of those.

In each of the two connected induced subgraphs above, repeatedly delete degree-one vertices.
The remaining two-cores are connected, disjoint, and have minimum degree at least two.
Deleting leaves preserves every cycle and the rank.
For large $r$, both cores therefore have rank at least $T/3$.
They are separated by the cut around $U_i$.
Apply Lemma~\ref{lem:random-protected-core-cut}.
Its normalized shores have ranks at least $T/3\ge r^{c/10}$ for sufficiently large $r$.
\end{proof}

\begin{lemma}[A quarter bound for any small cut with two large shores]
\label{lem:random-general-large-shore-quarter}
Fix positive constants $c,K,\delta$.
Suppose connected cubic graphs have rank $r\to\infty$, girth at least $c\log r$, and a cut of order $k\le K\log r$ whose shores are connected normalized poles of ranks at least $r^\delta$.
Then
\[
 c(G)\le\left(\frac14+o_{c,K,\delta}(1)\right)2^r.
\]
The cut is not required to have minimum cyclic order.
\end{lemma}

\begin{proof}
Write the shore ranks as $h,j$, so $r=h+j+k-1$.
If $k=1$, no cycle crosses the cut and $c(G)\le2^h+2^j=o(2^r)$.
Assume henceforth that $k\ge2$.
Theorem~\ref{thm:random-all-terminal-sectors-girth} applies to the first shore:
$k\le(\log h)^2$ for large $r$, and its required girth is $O_K(\log\log r)$.
Each even terminal sector therefore has at most $(1/2+o(1))2^h$ forests.
Lemma~\ref{lem:random-growing-port-aggregate} gives total forest mass at most
$(1/2+o(1))2^{j+k-1}$ in the other shore.
Every crossing cycle has a unique pair of forest restrictions, so its count is bounded by the product of these two quantities.
The internal cycle counts are at most $2^h+2^j=o(2^r)$.
All errors are uniform under the fixed constants in the statement.
\end{proof}

\begin{theorem}[A high-girth edge-count dichotomy]
\label{thm:random-high-girth-edge-dichotomy}
Fix $c>0$.
For all sufficiently large ranks $r$, let $G$ be a connected loopless cubic graph of girth at least $c\log r$.
Let $W$ be a shortest cycle and let $e=uv$ be any edge of $W$.
If
\[
 P_{uv}(G-e)>2^{r-2},
\]
then
\[
 c(G)<2^{r-1}.
\]
Equivalently, if $c(G)\ge2^{r-1}$, every edge of a shortest cycle belongs to at most $2^{r-2}$ cycles.
\end{theorem}

\begin{proof}
The cubic Moore bound gives $|V(W)|\le3\log r$ for large $r$.
The usual two-state injection, now using the existing edge $e$, gives
\[
 \frac{P_{uv}(G-e)}{2^{r-1}}
 \le\Pr(\xi_G-V(W)\text{ is a forest}).
\]
Explicitly, if $A=W-e$, a terminal path $Q$ gives the two distinct even states
$Q\cup\{e\}$ and $Q\mathbin\triangle A$.
Their restrictions outside $V(W)$ are forests.
The hypothesis therefore gives a marked forest probability greater than one half.

Lemma~\ref{lem:random-no-deletion-chain} supplies a cut of order at most $2A\log r$ with two polynomial-rank normalized shores.
Lemma~\ref{lem:random-general-large-shore-quarter} bounds the cycle count by
$(1/4+o(1))2^r$, which is strictly below $2^{r-1}$.
\end{proof}

\begin{corollary}[No logarithmic-girth minimal counterexample to the sharp order]
\label{cor:random-high-girth-minimal-counterexample}
Fix $c,K>0$.
Suppose the bound
\[
 c(H)\le2^{s-1}+K s2^{s/2}
\]
holds for every connected graph of rank $s=r-1$.
Then, for all sufficiently large $r$, every connected loopless cubic graph $G$ of rank $r$ and girth at least $c\log r$ satisfies
\[
 c(G)\le2^{r-1}+K r2^{r/2}.
\]
\end{corollary}

\begin{proof}
If $c(G)<2^{r-1}$, there is nothing to prove.
Otherwise choose an edge $e$ of a shortest cycle.
The theorem gives $P_e(G-e)\le2^{r-2}$.
The graph $G-e$ is connected and has rank $r-1$, so the assumed predecessor bound yields
\[
 c(G)=c(G-e)+P_e(G-e)
 \le2^{r-1}+K(r-1)2^{(r-1)/2}
 <2^{r-1}+Kr2^{r/2}.
\]
This argument requires the induction hypothesis for general connected graphs; it does not assume that $G-e$ is cubic.
\end{proof}

\begin{corollary}[The remaining girth window for a minimal counterexample]
\label{cor:random-minimal-counterexample-window}
There are absolute constants $K_*,C_3,C_4$ with the following property.
Fix $K\ge K_*$.
Suppose a connected loopless cubic graph $G$ of rank $r$ violates
\[
 c(G)\le2^{r-1}+Kr2^{r/2},
\]
while this bound holds for every connected graph of smaller rank.
If $r$ is sufficiently large, its girth $g$ and minimum cyclic cut order $k$ satisfy
\[
 \frac12\log\log r-\frac{11}{2}\log\log\log r-C_3
 <k=g
 <2\log\log r+2\log\log\log r+C_4.
\]
Moreover, every minimum cyclic cut has a shortest cycle as one shore.
The constants do not depend on $K$.
\end{corollary}

\begin{proof}
Choose $K_*$ at least the absolute constant in the established sharp-order bound for graphs with a cyclic cut of size at most four.
Such cuts are then excluded in $G$.
Corollary~\ref{cor:path-endpoint-loglog-cut} gives the stated lower bound on $k$, since $G$ has cycle count above the half bound.
In particular, $k$ and $g\ge k$ tend uniformly to infinity in the possible counterexamples.

Apply Corollary~\ref{cor:random-improved-girth-window}.
It excludes every girth between
\[
 \left\lceil2\log\log r+2\log\log\log r+C_*\right\rceil
 \quad\text{and}\quad
 \frac12\log r-\frac{11}{2}\log\log r-C_0.
\]
For sufficiently large $r$, the second endpoint exceeds $(\log r)/4$.
Corollary~\ref{cor:random-high-girth-minimal-counterexample}, with $c=1/4$, excludes every girth at least $(\log r)/4$.
The two ranges overlap.
Thus the upper bound on $g$ follows after increasing an absolute constant to account for the ceiling.

Finally, Corollary~\ref{cor:path-high-girth-noncycle-minimum-cut} excludes a minimum cyclic cut with two noncycle shores once $g$ is sufficiently large.
Every minimum cut therefore isolates a cycle.
Its boundary order equals that cycle's length, which is at least $g$.
Since $k\le g$, one has $k=g$, and the isolated cycle is shortest.
\end{proof}


\section{Checks and the remaining unrestricted question}
Each written argument has four independent full mathematical reviews.
The shortest-cycle identities also have exact finite checks on seventy-two random cubic graphs and a graph whose shortest five-cycle has two cyclic exterior components.
The checker verifies 292 prefix deletions, the terminal-path partitions, and independent cycle-space counts.
The cut-smoothing proof uses an explicit shore-partition lift and exact cycle bijections; it does not rely on numerical optimization.
An independent checker exhausts all cyclic vertex cuts of the Petersen and Heawood graphs and checks all 198 minimum-cut edge incidences.
The localization checker compares the complete restriction distributions in 1,193 cases, including 265 disconnected regions and 56 padded boundary cycles.
It enumerates 93,632 even states and verifies the rank formulas exactly.
The Moore geometry checker verifies 35,568 geodesics, 20,694 disjoint packed balls, and 10,579 port-free Moore-ball counts on 49 normalized-pole fixtures.
The affine-forest checker verifies translations and the exact saturation criterion on 247 fixtures and 8,312 affine states.

The unrestricted $O(r2^{r/2})$ bound is not proved.
A counterexample sequence after removal of the sharp four-cut class must have cyclic connectivity tending to infinity.
It must also have girth tending to infinity.
The finite recurrences and the quantitative localization theorem cover the displayed double-logarithmic range.
The remaining cyclic connectivity can still grow up to order $\log r$.
The high-girth affine-forest theorem forces each minimum cyclic cut in an above-half family to isolate a shortest cycle.
The no-deletion chain and the unconditional girth window reduce any sufficiently large minimal counterexample to the double-logarithmic range stated in the abstract.
For every fixed $C$, a remaining above-half family also has no cut of size at most $Cg$ with two connected shores whose ranks exceed the explicit constant $H_C$ in Section~\ref{sec:four-high-connectivity-weighted-kernel}.
A separate construction shows that an edge of a minimum cyclic cut need not lie on a short cycle: even a unique fixed-order minimum cut can have all its edges on cycles of length proportional to rank.
Thus a short-cycle selection argument would need an additional cycle-count hypothesis; minimum-cut geometry alone is insufficient.

The shared files \nolinkurl{CODEX-TO-FABLE-LIVE.txt} and \nolinkurl{RESEARCH-STATUS-2026-11-11.txt} record the current checkpoint and partial investigations.
The earlier formal modules are unchanged.
\end{document}
